% ==========================================
% THE MODULAR HOMOTOPY TYPE
% ==========================================

\begin{remark}[Type signature; modularity as a Level-5 reconstruction]
\label{rem:gc-type-signature}
\index{type signature!genus-complete chapter}
\index{modular reconstruction!hypothesis package}
Every theorem in this chapter carries the type signature
(\emph{Open quadrant}, \emph{factorization-algebra presentation on
$(X,D,\tau)$}, \emph{level $\bullet$}, \emph{hypothesis package}). The
genus tower lives across two Beilinson levels: the chain-level
modular functor and the curved $A_\infty$ structure are
\emph{level-$3$} statements (derived chiral centre as the
universal closed-sector cochain algebra acting on the chiral category
and its boundary modules); the scalar
genus class $F_g^{\mathrm{diag}}=\kappa(\cA)\lambda_g^{\mathrm{FP}}$
and the Hodge-class shadow on $\overline{\cM}_{g,n}$ are
\emph{level-$5$} statements obtained by trace + clutching on the
open factorization category. \emph{Modular invariance is the level-$5$
modular reconstruction}: it is not an intrinsic property of a closed
algebra, but the consequence, on the open factorization category, of
the trace + clutching package together with named analytic and
finiteness hypotheses (HS-sewing, weight-completed pro-conilpotency,
or Zhu finiteness with rational fusion when stronger forms are
invoked). Closed shadows have modular consequences; modularity itself
is recovered by reconstruction from the open lane. On the
\emph{KSDual} sublocus ($\Z/2$-fixed under $\cA\mapsto\cA^!$), the
modular reconstruction is an equivalence rather than an adjunction
with comparison map. Cross-volume specializations to
$\Phi_d^{(\Sigma_{d-1},C)}=\mathrm{Sp}^{\mathrm{ch}}_{\Sigma_{d-1},C}\circ
\Phi_d^{\mathrm{FA}}$ enter at level~$2$; the elliptic-fibre
specialization $C=E_\tau$ is the genus-$1$ shadow used in
\S\ref{rem:genus-complete-character} and below.
\end{remark}

The genus tower is governed by the modular Chern--Weil transform:
Maurer--Cartan data in the modular convolution algebra
$\mathfrak{g}_\cA^{\mathrm{mod}}$ are sent to tautological classes on
$\overline{\mathcal{M}}_{g,n}$. Its diagonal scalar trace is
\[
F_g^{\mathrm{diag}}(\cA)
=\kappa(\cA)\lambda_g^{\mathrm{FP}}.
\]
This is the whole free energy on the uniform-weight scalar lane and
at genus~$1$ for every modular Koszul algebra. Off that lane, the
all-weight free energy is
\[
F_g(\cA)=\kappa(\cA)\lambda_g^{\mathrm{FP}}
+\delta F_g^{\mathrm{cross}}(\cA),
\qquad
\delta F_1^{\mathrm{cross}}=0,
\]
with $\delta F_g^{\mathrm{cross}}$ supplied by mixed boundary graphs
(Theorem~\ref{thm:multi-weight-genus-expansion}). The higher
Chern--Weil projections carry the cubic, quartic, and all-degree
shadow obstruction tower; the all-degree MC element exists under the
hypotheses of Theorem~\ref{thm:mc2-bar-intrinsic}.

The moduli spaces $\overline{\mathcal{M}}_{g,n}$ form a modular
operad (Getzler--Kapranov~\cite{GeK98}): surfaces glue along
boundary divisors, and the gluing is associative up to coherent
homotopy. A chiral algebra $\cA$ on a curve $X$ is a coefficient
system for this operad. The \emph{Feynman transform}
$\mathrm{FT}_{\mathcal{M}\mathrm{od}}(\cA)$ (the bar complex
assembled over all genera) is therefore a single algebraic
object that \emph{presents} $\cA$ over the full modular operad,
exactly as the convolution algebra on the Steinberg variety
presents the Hecke algebra through the geometry of the
flag variety.

The result is a chain-level modular functor
(Theorem~\ref{thm:chain-modular-functor}): to each surface
$\Sigma_{g,n}$ a cochain complex, to each degeneration a chain
map, to each consistency relation a chain homotopy. Passing to
cohomology gives the classical cohomological modular-functor
package; for $\widehat{\fg}_k$ at integrable level,
$H^0$ is the space of conformal blocks, and the induced flat
connection has the KZ/Hitchin package as its integrable affine
comparison surface
(Remark~\ref{rem:moduli-variation}).

\emph{The diagonal scalar tower
$F_g^{\mathrm{diag}} = \kappa(\cA)\lambda_g^{\mathrm{FP}}$
\textup{(}Theorem~D\textup{)}
is the first Chern class of the determinant line bundle of the
bar family, computed by Grothendieck--Riemann--Roch.}  Mixed
weights contribute only through the separate cross-channel summand
$\delta F_g^{\mathrm{cross}}$.

\subsection{Genus one extensions}\label{subsec:elliptic}

\begin{definition}[Elliptic configuration spaces]\label{def:elliptic-config}
\index{configuration space!elliptic}
For a genus one curve $E_\tau$, let
$\overline{C}_n^{(1)}(E_\tau)$ be the Fulton--MacPherson
\textup{(}wonderful\textup{)} compactification of
$\operatorname{Conf}_n(E_\tau)$. It is obtained by iteratively blowing
up the proper transforms of the partial diagonals
\[
D_I=\{z_i=z_j\text{ for all }i,j\in I\}\subset E_\tau^n,
\qquad |I|\geq2,
\]
in an order compatible with inclusion. The divisor over~$D_I$ maps to
$D_I\cong E_\tau^{\,n-|I|+1}$ and has normal fibre
\[
\mathbb P\bigl(T_{E_\tau}^{I}/T_{E_\tau}\bigr)
\cong \mathbb P^{|I|-2},
\]
because $T_{E_\tau}$ is trivial. For $|I|=2$ the diagonal already has
codimension~$1$; the boundary record is the logarithmic normal
direction, not a new exceptional elliptic curve.
\end{definition}

\begin{theorem}[Elliptic bar complex; \ClaimStatusProvedElsewhere{} \cite{FBZ04}]\label{thm:elliptic-bar}
\textbf{Type signature.} (Open quadrant; factorization-algebra
presentation on $(E_\tau, D, \tau)$; level~$2$ for the chain-level bar
construction; hypothesis package: chiral algebra $\cA$ with
finite-type or weight-completed bar coalgebra $B(\cA)$.)
For a chiral algebra $\mathcal{A}$ on $E_\tau$, the elliptic bar complex is:
\[\bar{B}^{(1),n}_{\text{geom}}(\mathcal{A}) = \Gamma\left(\overline{C}_{n+1}^{(1)}(E_\tau), j_*j^*\mathcal{A}^{\boxtimes(n+1)} \otimes \Omega^n_{ell}(\log D)\right)\]
where $\Omega^n_{ell}$ consists of meromorphic forms with logarithmic poles, satisfying:
\begin{enumerate}
\item \emph{Elliptic periodicity:} Forms are doubly periodic modulo residues.
\item \emph{Modular covariance of the chain-level functor.} Under
$\tau \mapsto (a\tau+b)/(c\tau+d)$, the elliptic bar complex transforms
covariantly with weight determined by the conformal weights of the
generators. This is a level-$5$ statement on the open factorization
category: the modular covariance is the $\mathrm{SL}(2,\Z)$-equivariance
of the trace + clutching map on $\overline{\cM}_{1,n}$, not an intrinsic
property of $\cA$ as a closed algebra.
\item \emph{Theta function expansion:} Near divisors, forms expand in
Jacobi theta functions.
\end{enumerate}

The differential decomposes as:
\[d^{(1)} = d_{\text{local}} + d_{\text{global}} + d_{\text{quantum}}\]
where $d_{\text{quantum}}$ encodes the holonomy around non-contractible cycles.
\end{theorem}

\begin{proof}[Construction sketch]
Elliptic curves are group varieties, so configuration spaces inherit a group action. This forces specific structures:

\emph{Local structure.} Near collisions, the OPE universality theorem ensures local behavior matches genus zero. In coordinates $(u, \epsilon_{ij}, \theta_{ij})$ near $D_{ij}$:
\[\eta_{ij}^{(1)} = d\log\epsilon_{ij} + id\theta_{ij} + \text{elliptic corrections}\]

\emph{Global elliptic functions.} By Liouville's theorem, meromorphic doubly-periodic functions satisfy:
\[\sum_{\text{poles in } \mathcal{F}} \text{Res}_p[f] = 0\]
where $\mathcal{F}$ is a fundamental domain. This constraint modifies the bar differential.

\emph{Theta function basis.} The Jacobi theta functions provide the natural basis:
\begin{align}
\vartheta_1(z|\tau) &= -i\sum_{n \in \mathbb{Z}}(-1)^n q^{(n+1/2)^2/2} e^{2\pi i(n+1/2)z} \\
\vartheta_2(z|\tau) &= \sum_{n \in \mathbb{Z}} q^{(n+1/2)^2/2} e^{2\pi i(n+1/2)z} \\
\vartheta_3(z|\tau) &= \sum_{n \in \mathbb{Z}} q^{n^2/2} e^{2\pi inz} \\
\vartheta_4(z|\tau) &= \sum_{n \in \mathbb{Z}}(-1)^n q^{n^2/2} e^{2\pi inz}
\end{align}
where $q = e^{2\pi i\tau}$.

\emph{Quantum differential.} The genus-$1$ fiberwise differential $\dfib$ (Convention~\ref{conv:higher-genus-differentials}) acquires monodromy contributions from the two cycles of $E_\tau$. With $\omega = dz$ the normalized holomorphic differential ($\oint_A \omega = 1$, $\oint_B \omega = \tau$), the period matrix $\Omega_1 = \tau$ enters through the Gauss--Manin connection: sections of the bar complex acquire monodromy $M_A = \mathrm{id}$, $M_B = e^{2\pi i \tau \cdot \partial_z}$ around the $A$- and $B$-cycles respectively.

\emph{Modular transformation.} Under $SL_2(\mathbb{Z})$:
\[\vartheta_1\left(\frac{z}{c\tau+d}\bigg|\frac{a\tau+b}{c\tau+d}\right) = \epsilon(a,b,c,d)\sqrt{c\tau+d}\,e^{\frac{\pi i cz^2}{c\tau+d}}\vartheta_1(z|\tau)\]
where $\epsilon$ is an 8th root of unity determined by $(a,b,c,d) \mod 2$.

\emph{Elliptic distance formula.} The proper distance on the elliptic curve:
\[|z_i - z_j|_{E_\tau} = |\sigma(z_i - z_j; \tau)|e^{-\eta(\tau)\text{Im}(z_i - z_j)^2/\text{Im}(\tau)}\]
where $\sigma$ is the Weierstrass sigma function and $\eta$ is the Weierstrass eta value $\eta_1 = \zeta(\omega_1)$ (not the Dedekind eta function).
\end{proof}

\begin{example}[Free fermion at genus one]
Consider the free fermion chiral algebra with generators $\psi(z), \psi^*(z)$ satisfying:
\[\psi(z)\psi^*(w) \sim \frac{1}{z-w}\]

At genus one with spin structure $\alpha \in \mathbb{Z}_2 \times \mathbb{Z}_2$, we have four sectors:

\begin{center}
\begin{tabular}{|c|c|c|}
\hline
Sector & Boundary Conditions & Partition Function \\
\hline
NS-NS & $\psi(z+1) = -\psi(z), \psi(z+\tau) = -\psi(z)$ & $\frac{\vartheta_3(0|\tau)}{\eta(\tau)}$ \\
NS-R & $\psi(z+1) = -\psi(z), \psi(z+\tau) = \psi(z)$ & $\frac{\vartheta_4(0|\tau)}{\eta(\tau)}$ \\
R-NS & $\psi(z+1) = \psi(z), \psi(z+\tau) = -\psi(z)$ & $\frac{\vartheta_2(0|\tau)}{\eta(\tau)}$ \\
R-R & $\psi(z+1) = \psi(z), \psi(z+\tau) = \psi(z)$ & $\frac{\vartheta_1(0|\tau)}{\eta(\tau)} = 0$ \\
\hline
\end{tabular}
\end{center}

The bar complex calculation:
\[Z_\alpha = \int_{\overline{C}_*^{(1)}(E_\tau)} \exp\left(\sum_{i<j} \log\frac{\vartheta_\alpha(z_i - z_j|\tau)}{\vartheta_1(z_i - z_j|\tau)} \cdot \eta_{ij}^{(1)}\right)\]

The R-R sector vanishes: $\vartheta_1(0|\tau) = 0$ identically (the odd theta function vanishes at the origin), so $Z_{\mathrm{R\text{-}R}} = 0$. The other three sectors are non-vanishing.
\end{example}

\begin{theorem}[Extension obstruction; \ClaimStatusConditional{} \cite{FBZ04,CG17}]\label{thm:extension-obstruction}
\index{extension obstruction!genus one}
\textbf{Type signature.} (Open quadrant; factorization-algebra
presentation on $(E_\tau, D, \tau)$; level~$5$, modular reconstruction
on the genus-$1$ open category; hypothesis package: $\cA_0$ a
genus-$0$ conformal chiral algebra; FBZ trace + clutching package
satisfied; bar coalgebra finite-type or weight-completed.)

A genus-zero conformal chiral algebra $\mathcal{A}_0$ whose torus
sewing data satisfy the FBZ trace + clutching hypotheses
\textup{(}see~\cite{FBZ04}\textup{)} extends to genus~$1$ as a projective
factorization algebra on $\overline{\cM}_{1,n}$, twisted by a power of
the Hodge line. The obstruction to a strict untwisted extension is
the genus-$1$ scalar curvature class
\[
\mathrm{Obs}_1(\mathcal{A}_0)
=\kappa(\mathcal{A}_0)\lambda_1
\in H^2(\overline{\mathcal{M}}_{1,n},\mathcal{L}_{\mathcal{A}}).
\]
For the Virasoro family, $\kappa(\mathrm{Vir}_c)=c/2$, hence
$\mathrm{Obs}_1(\mathrm{Vir}_c)=(c/2)\lambda_1$; it vanishes exactly
at $c=0$. The integrated Virasoro obstruction is
$(c/2)\int_{\overline{\mathcal{M}}_{1,1}}\lambda_1=c/48$ on the
stack, or $c/24$ on the coarse moduli space. When
$c\in24\mathbb Z$, the character phase $e^{-2\pi i c/24}$ is
trivial, but the factorization algebra is still naturally twisted by
$\mathbb{E}^{c/2}$. The hypothesis package decomposes as:
\begin{enumerate}
\item \emph{Level-$1$ datum (chart curvature).} The chart curvature
$\kappa(\mathcal{A}_0)$ is read off the OPE; for Virasoro,
$\kappa(\mathrm{Vir}_c)=c/2$.
\item \emph{Level-$5$ datum (modular reconstruction; trace + clutching).}
The genus-$1$ characters $\chi_i(\tau)$ assemble into a vector-valued
modular form precisely because the trace + clutching package for the
open category lifts to the $\mathrm{SL}(2,\Z)$-equivariant
$\overline{\cM}_{1,1}$-extension. \emph{Modularity is not an intrinsic
closed-algebra property: the closed shadow has modular consequences,
but those consequences arise from the open-category trace and
clutching, not from a closure construction.} On the \emph{KSDual}
sublocus this reconstruction becomes an equivalence; off that
sublocus, additional finiteness data
\textup{(}\textsc{C}$_2$-cofinite/Zhu-finite, rational fusion, or HS-sewing
\textup{(}Definition~\ref{def:hs-sewing}\textup{)}, or the
weight-completed pro-conilpotent ambient\textup{)} enter as named hypotheses.
\item \emph{Level-$5$ datum (integrality).} Fusion rules
$N_{ij}^k \in \Z_{\ge 0}$ are part of the same trace + clutching
package on the open category; they are not closed-algebra integrality
constraints in disguise.
\end{enumerate}

Here $\lambda_1=c_1(\mathbb{E})$ is the first Chern class of the
Hodge bundle. This is the genus-$1$ specialization of the scalar
diagonal formula of Theorem~\ref{thm:genus-universality}; no
higher-genus cross-channel term appears at $g=1$.
\end{theorem}

% ==========================================
% HIGHER GENUS THEORY (g ≥ 2)
% ==========================================

\subsection{Higher genus theory}\label{subsec:higher-genus}

\begin{definition}[Higher-genus configuration spaces]\label{def:higher-genus-config}
\index{configuration space!higher genus}
For a genus $g \geq 2$ curve $\Sigma_g$, the configuration space has a fibration structure:
\[\pi: \overline{C}_n^{(g)}(\Sigma_g) \to \overline{\mathcal{M}}_{g,n}\]
with fiber $(\Sigma_g)^n_{\text{ordered}}$ over each point in moduli space.

The boundary stratification has three types: collision divisors $D_{ij}^{(g)}$ ($z_i \to z_j$), separating divisors $D_{I|J}^{\text{sep}}$ ($\Sigma_g \to \Sigma_{g_1} \cup \Sigma_{g_2}$, $g_1 + g_2 = g$), and non-separating divisors $D_\gamma^{\text{nonsep}}$ (cycle $\gamma$ pinched). Each stratum carries a natural orientation from the complex structure.
\end{definition}

\begin{theorem}[Higher-genus bar differential; \ClaimStatusConditional{} \cite{CG17,costello-renormalization}]\label{thm:higher-genus-diff}
\index{bar differential!higher genus}
\textbf{Type signature.} (Open quadrant; bar/twisting presentation on
$(\Sigma_g, D, \tau)$; level~$2$ for the chain-level bar differential,
level~$3$ for the curved $A_\infty$ structure on the derived chiral
centre; hypothesis package: chiral algebra $\cA$ on $\Sigma_g$ with
finite-type or weight-completed bar coalgebra, Leray degeneration on
the universal curve, Koszul weight filtration.)
The total corrected bar differential at genus $g$
\textup{(}Convention~\textup{\ref{conv:higher-genus-differentials})} has the explicit form:
\[\Dg{g} = d_{\text{local}} + d_{\text{period}} + d_{\text{moduli}} + d_{\text{quantum}}\]

where:
\begin{align}
d_{\text{local}} &= \sum_{i<j} \text{Res}_{D_{ij}} \left[\mu_{ij} \otimes \eta_{ij}^{(g)}\right] \\
d_{\text{period}} &= \sum_{k=1}^{2g} \oint_{C_k} \eta_k \cdot \delta_{C_k^*} \\
d_{\text{moduli}} &= \sum_{a=1}^{3g-3} \frac{\partial}{\partial \tau_a} \otimes d\tau_a \\
d_{\text{quantum}} &= \sum_{\Gamma \in \mathcal{G}_{g,n}} \frac{1}{|\text{Aut}(\Gamma)|} \int_{\mathcal{M}_\Gamma} \omega_\Gamma
\end{align}

Here:
\begin{itemize}
\item $C_k$ ($k = 1, \ldots, 2g$) are the homology cycles ($A_1,\ldots,A_g,B_1,\ldots,B_g$), $C_k^*$ their Poincar\'e dual cohomology classes, and $\eta_k$ are the harmonic representatives of $H^1_{\mathrm{dR}}(\Sigma_g)$ dual to $C_k$
\item $\tau_a$ are Teichmüller coordinates on $\mathcal{M}_g$
\item $\mathcal{G}_{g,n}$ is the set of stable graphs of genus $g$ with $n$ legs
\item $\omega_\Gamma$ is the form associated to graph $\Gamma$
\end{itemize}
\end{theorem}

\begin{proof}[Construction via relative cohomology]
Consider the universal curve $\pi: \mathcal{C}_g \to \mathcal{M}_g$. The derived pushforward is computed via the bar complex.
Since the fibers of $\pi$ are smooth curves, $R^n\pi_*\mathcal{A} = 0$ for $n \geq 2$, so $R\pi_*\mathcal{A}$ is an object of cohomological amplitude $[0,1]$. When the Leray spectral sequence for~$\pi$ with coefficients in~$\mathcal{A}$ degenerates at~$E_2$ and the resulting filtration splits, one obtains:
\[R\pi_* \left(\mathcal{A}|_{\mathcal{C}_g}\right) \simeq R^0\pi_* \mathcal{A} \oplus R^1\pi_* \mathcal{A}[-1].\] For constant coefficients, this follows from Deligne's theorem on smooth proper morphisms; for general chiral algebra coefficients, the degeneration is an additional hypothesis that holds for Koszul algebras by the weight filtration on the associated graded bar complex (Definition~\ref{def:bar-cobar-filtration}).

By Grothendieck's theorem, the Hodge--Leray spectral sequence reads:
\[E_2^{p,q} = H^p(\mathcal{M}_g, R^q\pi_*\mathcal{A}) \Rightarrow H^{p+q}(\bar{B}^{(g)}(\mathcal{A}))\]

Each piece of the $E_2$ page contributes a component of the differential:
\begin{itemize}
\item $E_2^{0,*}$: Local OPE contributions ($d_{\text{local}}$)
\item $E_2^{1,*}$: Period-dependent corrections from $H^1(\mathcal{M}_g, R^0\pi_*\mathcal{A})$ ($d_{\text{period}}$). (This is \emph{sheaf} cohomology of $R^0\pi_*\mathcal{A}$ on $\mathcal{M}_g$, which can be nontrivial even though $H^1(\mathcal{M}_g, \mathbb{Q}) = 0$ for $g \geq 2$ by Harer's theorem; the period parameters come from $H^1(\Sigma_g, \mathbb{C})$ via the local system $R^1\pi_*\underline{\mathbb{C}}$.)
\item $E_2^{p,*}$ ($p \geq 2$): Higher moduli contributions ($d_{\text{moduli}}$)
\end{itemize}
The quantum differential $d_{\text{quantum}}$ arises from the curved $A_\infty$ structure at genus~$g$ (the curving element $m_0^{(g)}$ depending on the moduli), which is encoded in the extension data of the spectral sequence (the $E_2$ page alone determines the graded pieces, but the extensions encode the quantum corrections).

The identity $\Dg{g}^{\,2} = 0$ for the total differential at the convolution level is equivalent to the Maurer--Cartan equation for the curved $A_\infty$ structure, which is proved by genus induction in Theorem~\ref{thm:genus-induction-strict}. (At the ambient level with planted-forest corrections, $D^2 = 0$ is Theorem~\ref{thm:ambient-d-squared-zero}.)
Concretely, the curving element $m_0^{(g)}$ at each genus is the
genus-$g$ component of the universal Maurer--Cartan class
$\Theta_\cA$
(Theorem~\ref{thm:mc2-bar-intrinsic}). On the uniform-weight
scalar-diagonal lane,
$\Theta_\cA^{\min}=\kappa(\cA)\eta\otimes\Lambda$
(Theorem~\ref{thm:explicit-theta}), and the scalar component is
$m_{0,\mathrm{diag}}^{(g)}=\kappa(\cA)\lambda_g$, with
$\lambda_g=c_g(\mathbb{E})$ on $\overline{\mathcal{M}}_g$.
For multi-weight algebras at $g\geq2$, the diagonal equality
$\Gamma_\cA=\kappa(\cA)\Lambda$ does not determine the full free
energy: mixed generator channels contribute
$\delta F_g^{\mathrm{cross}}(\cA)$
(Theorem~\ref{thm:multi-weight-genus-expansion}). The full MC equation
$\ell_1(\Theta) + \tfrac{1}{2}\ell_2(\Theta,\Theta) = 0$ assembles
the genus-by-genus nilpotency $\sum_{g_1+g_2=g} d_{g_1} \circ d_{g_2}
= 0$ into a single algebraic identity, with $\Theta$ encoding both
the period-dependent corrections (from sheaf cohomology of
$R^0\pi_*\cA$) and the quantum corrections (from higher moduli
contributions).
\end{proof}

\begin{example}[WZW model at higher genus]
\index{WZW model!higher genus}
For the $\widehat{\mathfrak{g}}_k$ WZW model on $\Sigma_g$, the partition function is given by the Verlinde formula:
\begin{equation}\label{eq:verlinde-general-wzw}
Z_g(k) = \sum_{\lambda \in \hat{P}_+^k} \left(\frac{S_{0\lambda}}{S_{00}}\right)^{2-2g}
\end{equation}
where $\hat{P}_+^k$ denotes the set of level-$k$ integrable highest weights, and $S$ is the modular $S$-matrix.

The bar complex at genus~$g$ uses the Arakelov--Green function on~$\Sigma_g$ as its propagator:
\[G_g(z,w) = -\log|E(z,w)|^2 + 2\pi\sum_{j,\ell=1}^g \mathrm{Im}\!\int_z^w \omega_j \cdot (\mathrm{Im}\,\Omega)^{-1}_{j\ell} \cdot \mathrm{Im}\!\int_z^w \omega_\ell\]
where $E(z,w)$ is the prime form (Appendix~\ref{app:theta}), $\omega_1, \ldots, \omega_g$ are the normalized holomorphic differentials ($\oint_{A_j}\omega_k = \delta_{jk}$), and $\Omega$ is the period matrix. The second term subtracts the zero-mode contribution, ensuring $G_g$ is well-defined on~$\Sigma_g$ (which admits no global Green's function without this regularization).

The bar complex correlation functions (Theorem~\ref{thm:mk-general-structure-vol1}) are computed by summing over stable graphs with edges weighted by~$G_g$ and vertices weighted by OPE structure constants of~$\widehat{\mathfrak{g}}_k$. The resulting genus expansion recovers~\eqref{eq:verlinde-general-wzw}: the exponent $(2-2g)$ arises from the Euler characteristic $\chi(\Sigma_g) = 2 - 2g$ appearing in the genus-$g$ propagator count.
\end{example}

% ==========================================
% UNIVERSAL TOWER AND SPECTRAL SEQUENCES
% ==========================================

\subsection{The universal genus tower}\label{subsec:universal-tower}
\index{genus tower}

\begin{theorem}[Master tower of boundary residues; \ClaimStatusProvedHere]\label{thm:master-tower}
\textbf{Type signature.} (Open quadrant; modular-operad/Feynman-transform
presentation on $\{\overline{\cM}_{g,n}\}$; level~$2$ for the bar
boundary residues, level~$3$ for the assembled Feynman-transform bulk;
hypothesis package: chiral algebra $\cA$ with Koszul/finite-type or
weight-completed bar coalgebra at every fixed window $(G,R)$;
Getzler--Kapranov modular operad structure on
$\{\overline{\cM}_{g,n}\}_{g,n}$.)
The genus tower is a boundary-residue tower, not a chain of
forgetful maps $\bar B^{(g+1)}\to\bar B^{(g)}$.  For each stable
pair $(g,n)$ the boundary divisors of $\overline{\mathcal{M}}_{g,n}$
define residue maps
\begin{itemize}
\item \emph{Nonseparating degeneration.}
\[
\rho_{\mathrm{irr}}\colon
\bar{B}^{(g)}_n(\mathcal{A})
\longrightarrow
\mathrm{Tr}_{\mathcal{A}}\bigl(\bar{B}^{(g-1)}_{n+2}(\mathcal{A})\bigr),
\qquad
\alpha\longmapsto \operatorname{Res}_{D_{\mathrm{irr}}}(\alpha),
\]
where $\mathrm{Tr}_{\mathcal{A}}$ is the self-sewing contraction over
the two new markings.
\item \emph{Separating degeneration.}
\[
\rho_{g_1,I}\colon
\bar{B}^{(g)}_n(\mathcal{A})
\longrightarrow
\bar{B}^{(g_1)}_{|I|+1}(\mathcal{A})
\,\widehat{\otimes}\,
\bar{B}^{(g_2)}_{|I^c|+1}(\mathcal{A}),
\qquad
\alpha\longmapsto \operatorname{Res}_{D_{g_1,I}}(\alpha),
\]
with $g_1+g_2=g$.
\end{itemize}
These are the cooperadic structure maps dual to the
Getzler--Kapranov clutching maps.

The total complex with string coupling $g_s$:
\[
\bar{B}^{\mathrm{total}}(\mathcal{A})
=
\varprojlim_{G,R}
\bigoplus_{\substack{0\leq g\leq G\\0\leq r\leq R}}
g_s^{2g-2}\,\bar{B}^{(g),\leq r}(\mathcal{A}).
\]
The completion is genus-adic and bar-length filtered: every finite
window $(G,R)$ is an ordinary finite direct sum, and the completed
object is the inverse limit of these windows.

The levels $g=0,1,\geq 2$ correspond to tree-level (classical), one-loop, and higher-loop amplitudes respectively.
\end{theorem}

\begin{proof}
The modular operad structure on
$\{\overline{\mathcal{M}}_{g,n}\}_{g,n}$ (Getzler--Kapranov)
provides clutching and self-gluing maps along boundary divisors.
The Feynman transform identification
$\bar{B}^{\mathrm{full}}(\mathcal{A}) \simeq
\mathrm{FT}_{\mathcal{M}\mathrm{od}}(\mathcal{A})$
(Theorem~\ref{thm:prism-higher-genus}(ii)) assembles the
genus-$g$ bar complexes into a modular coalgebra. The maps
$\rho_{\mathrm{irr}}$ and $\rho_{g_1,I}$ are the residue maps to
the irreducible and separating boundary strata.  They are compatible
with the bar differential because the Feynman transform differential
is the signed boundary operator of the modular operad; the two
codimension-$2$ presentations of a boundary stratum cancel with
opposite signs.
The finite-window formula above is stable under these maps: residues
do not increase genus and the bar differential raises or preserves the
bar-length filtration by a bounded amount on each window. Taking the
inverse limit gives the completed genus tower and its continuous
differential.
\end{proof}

\subsection{Chain-level modular functor}\label{subsec:chain-modular-functor}
\index{modular functor!chain-level}

The diagonal scalar genus tower
$F_g^{\mathrm{diag}}(\cA)=\kappa(\cA)\lambda_g^{\mathrm{FP}}$
is the first Chern class of a chain-level modular object.
For multi-weight algebras the full free energy also contains
$\delta F_g^{\mathrm{cross}}$ (Theorem~\ref{thm:multi-weight-genus-expansion});
the chain-level object records the data before this scalar projection.
The bar complex refines the classical modular-functor
formalism (Segal, Bakalov--Kirillov) to cochain complexes with
$A_\infty$-morphisms as structure maps, and it is this
chain-level modular functor (not the scalar genus tower) that
constitutes the full modular homotopy invariant of a chiral
algebra (Remark~\ref{rem:four-levels}).
The hierarchy is:
\begin{itemize}
\item \emph{Level~$0$:}
 $\kappa(\cA) \cdot \hat{A}$-genus: numbers
 (Theorem~\ref{thm:modular-characteristic}).
\item \emph{Level~$1$:}
 $H^*(\barB^{(g)}(\cA), \Dg{g})$: graded vector spaces
 (the derived fibers of the factorization homology bundle).
 For $\widehat{\fg}_k$ at integrable level, $H^0$ recovers
 the Tsuchiya--Ueno--Yamada space of conformal blocks, and
 its dimension is the Verlinde number
 (Example~\ref{eq:verlinde-general-wzw}, Remark~\ref{rem:chain-vs-classical-mf}).
\item \emph{Level~$2$:}
 Variation of $H^*$ over $\mathcal{M}_g$: a flat connection.
 On the integrable affine lane, its monodromy furnishes the
 bar-side comparison surface for the KZ/Hitchin package and
 the resulting mapping class group representation on conformal blocks.
\item \emph{Level~$3$:}
 Factorization and sewing: the following theorem.
\end{itemize}

\begin{theorem}[Chain-level boundary-residue modular functor from bar complex;
\ClaimStatusProvedHere]\label{thm:chain-modular-functor}
\index{modular functor!from bar complex}
\textbf{Type signature.} (Open quadrant; bar/Feynman-transform
presentation; level~$3$ for the chain assignment, level~$5$ for the
modular-functor axiomatic; hypothesis package: Koszul chiral algebra
$\cA$ with finite-type bar coalgebra, plus the trace + clutching
package on $\{\overline{\cM}_{g,n}\}$. The functor's mapping-class
$S$ and $T$ data are the modular reconstruction of $\cA$ from its
open-category trace, not an intrinsic closed-algebra modularity.)
For a Koszul chiral algebra $\mathcal{A}$, the finite
genus--bar-length windows of
$\bar{B}^{\mathrm{total}}(\mathcal{A})$ from
Theorem~\ref{thm:master-tower} satisfy the axioms of a
homotopy-coherent boundary-residue modular functor; the inverse limit
is the completed continuous object:
\begin{enumerate}[label=\textup{(\roman*)}]
\item \emph{Assignment}: To each genus~$g$ with $n$ marked points,
 a cochain complex $V_{g,n} := \bar{B}^{(g)}_n(\mathcal{A})$
 with $d_{g,n}^2 = 0$
 in the square-zero total-differential model
 \textup{(}Theorem~\textup{\ref{thm:genus-induction-strict})}; in the
 genuinely curved model the same datum is read in the coderived
 category, with curvature $m_0^{(g)}$.
\item \emph{Factorization}: For each separating degeneration
 $g = g_1 + g_2$, a boundary-residue chain map
 \begin{equation}\label{eq:factorization-map}
 \Delta_{\mathrm{sep}}: V_{g,n} \to
 \bigoplus_{\substack{g_1+g_2=g \\ I \sqcup J = \{1,\ldots,n\}}}
 V_{g_1,|I|+1} \otimes V_{g_2,|J|+1}
 \end{equation}
 induced by residues along the separating boundary divisor
 $D_{\mathrm{sep}} \subset \overline{\mathcal{M}}_{g,n}$.
\item \emph{Self-sewing}: For each nonseparating degeneration, a boundary-residue chain map
 \begin{equation}\label{eq:self-sewing-map}
 \Delta_{\mathrm{nonsep}}: V_{g,n} \to V_{g-1,n+2}
 \end{equation}
 induced by residues along the nonseparating boundary divisor
 $D_{\mathrm{nonsep}} \subset \overline{\mathcal{M}}_{g,n}$.
\item \emph{Coherence}: The maps $\Delta_{\mathrm{sep}}$ and
 $\Delta_{\mathrm{nonsep}}$ satisfy the Segal--Moore--Seiberg
 consistency relations up to explicit chain homotopies
 provided by the $A_\infty$ operations $m_k$ ($k \geq 3$).
\end{enumerate}
\end{theorem}

\begin{proof}
Part~(i) is Theorem~\ref{thm:genus-induction-strict}.

Parts~(ii)--(iii): The Feynman transform identification
$\bar{B}^{\mathrm{full}}(\mathcal{A}) \simeq
\mathrm{FT}_{\mathcal{M}\mathrm{od}}(\mathcal{A})$
(Theorem~\ref{thm:prism-higher-genus}(ii)) equips the genus-graded
bar complex with the Feynman-transform structure associated to the
modular operad $\{\overline{\mathcal{M}}_{g,n}\}_{g,n}$
(Getzler--Kapranov \cite{GeK98}). Applying boundary residues gives
the corresponding modular-coalgebra maps:
\begin{itemize}
\item \emph{Separating boundaries}: Residues along
 $D_{g_1,I} \cong \overline{\mathcal{M}}_{g_1,|I|+1}
 \times \overline{\mathcal{M}}_{g_2,|J|+1}$ give the
 coproduct~\eqref{eq:factorization-map}.
\item \emph{Nonseparating boundaries}: Residues along
 $D_\gamma \cong \overline{\mathcal{M}}_{g-1,n+2}$ give the
 self-sewing map~\eqref{eq:self-sewing-map}.
\end{itemize}
That these are chain maps (commute with the bar differential)
follows from the compatibility of the Feynman transform differential
with the modular operad structure \cite[Theorem~5.4]{GeK98}.

Part~(iv): The consistency relations (associativity of
iterated factorization, commutativity of factorization with
self-sewing) hold at the level of boundary strata of
$\overline{\mathcal{M}}_{g,n}$: the codimension-$2$ boundary
$D_{\mathrm{sep},1} \cap D_{\mathrm{sep},2}$ decomposes in two
ways, yielding the Moore--Seiberg consistency condition. At chain
level, this equality holds only up to homotopy: the bar complex
carries a curved $A_\infty$ structure
(Theorem~\ref{thm:bar-ainfty-complete}), and the higher
operations $m_k$ ($k \geq 3$) provide the homotopies
$\Delta_1 \circ \Delta_2 \simeq \Delta_2 \circ \Delta_1$
relating different orderings of degeneration. The homotopy
coherence data is encoded in the $A_\infty$ relations
$\sum_{r+s+t=n} (-1)^{rs+t} m_{r+1+t} \circ
(\mathrm{id}^{\otimes r} \otimes m_s \otimes
\mathrm{id}^{\otimes t}) = 0$.
\end{proof}

\begin{remark}[Comparison with classical modular functors]
\label{rem:chain-vs-classical-mf}
\index{modular functor!classical vs chain-level}
Passing to $H^*(V_{g,n})$ gives the classical cohomological modular-functor package~\cite{BK01}; for $\widehat{\mathfrak{g}}_k$ at integrable level, $H^0(V_{g,n})$ is the Tsuchiya--Ueno--Yamada space of conformal blocks. The chain-level refinement is necessary at generic~$k$ (multi-degree bar complex) and when $\kappa(\cA) \neq 0$ (curvature $m_0^{(g)} = \kappa(\cA) \cdot \lambda_g$ ; Theorem~\ref{thm:genus-graded-bar}).
\end{remark}

\begin{corollary}[Verdier-dual modular functors;
\ClaimStatusProvedHere]\label{cor:dual-modular-functor}
\index{modular functor!Koszul dual}
\textbf{Type signature.} (Open quadrant; Verdier-Koszul reflection at
level~$2$, with the chain-level modular functor at level~$3$/$5$;
hypothesis package: finite-type Koszul dual coalgebra $\cA^i$;
genuswise Poincar\'e--Koszul pairing perfect; Verdier duality on
$\overline{\cM}_{g,n}$. On the \emph{KSDual} sublocus
\textup{(}$\Z/2$-fixed under $\cA\mapsto\cA^!$\textup{)} the dualised
functor agrees with the original; off that sublocus the equality is
the Verdier-Koszul reflection with comparison map.)
Koszul dual coalgebra and, under finite-type dualizability,
$\mathcal{A}^!:=(\mathcal{A}^i)^\vee$ the Koszul dual algebra. Assume
the genuswise Poincar\'e--Koszul pairing of
Theorem~\ref{thm:poincare-extended} is perfect. The chain-level
modular functor attached to $\mathcal{A}^!$ is the Verdier dual of the
bar functor attached to $\mathcal{A}$:
\begin{equation}\label{eq:dual-mf}
V_{g,n}(\mathcal{A}^!) \simeq
\mathbb D_{\overline{\mathcal{M}}_{g,n}}
\bigl(V_{g,n}(\mathcal{A})\bigr),
\qquad
\mathbb D_M(-):=\mathrm{RHom}(-,\omega_M^\bullet).
\end{equation}
The factorization and self-sewing maps are interchanged by duality:
$\Delta_{\mathrm{sep}}(\mathcal{A}^!) = \Delta_{\mathrm{sep}}
(\mathcal{A})^\vee$ and $\Delta_{\mathrm{nonsep}}(\mathcal{A}^!)
= \Delta_{\mathrm{nonsep}}(\mathcal{A})^\vee$.
\end{corollary}

\begin{proof}
Theorem~\ref{thm:poincare-extended} gives the Verdier duality
identification.  Compatibility with factorization and self-sewing
maps follows from naturality of Verdier duality for the boundary
embeddings
$\overline{\mathcal{M}}_{g_1,|I|+1}\times
\overline{\mathcal{M}}_{g_2,|I^c|+1}\hookrightarrow
\overline{\mathcal{M}}_{g,n}$ and
$\overline{\mathcal{M}}_{g-1,n+2}\hookrightarrow
\overline{\mathcal{M}}_{g,n}$.
\end{proof}

\begin{remark}[Moduli variation and the integrable affine KZ/Hitchin comparison]%
\label{rem:moduli-variation}%
\index{KZ connection!from bar complex}%
\index{Hitchin connection!from bar complex}%
The chain-level modular functor (Theorem~\ref{thm:chain-modular-functor})
produces, at each genus~$g$, a sheaf of cochain complexes
$V_{g,n}$ over $\overline{\mathcal{M}}_{g,n}$.
Its cohomology $\mathcal{V}_{g,n} := H^*(V_{g,n})$
is a coherent sheaf on $\overline{\mathcal{M}}_{g,n}$
(a local system on the interior $\mathcal{M}_{g,n}$), and
the Gauss--Manin connection arising from the variation of
the bar differential over moduli space equips it with
a \emph{flat connection}:
\[
\nabla^{\mathrm{GM}}\colon
\mathcal{V}_{g,n} \to
\mathcal{V}_{g,n} \otimes \Omega^1_{\mathcal{M}_{g,n}}.
\]

For $\widehat{\mathfrak{g}}_k$ at integrable level, this flat
connection furnishes the bar-side comparison surface for the
\emph{KZ/Hitchin package}
on conformal blocks
(Tsuchiya--Ueno--Yamada~\cite{TUY89};
Hitchin~\cite{Hitchin90}).
Its monodromy gives a representation
$\pi_1(\mathcal{M}_{g,n}) \to \mathrm{GL}(\mathcal{V}_{g,n})$
of the mapping class group, the genus-$g$
modular representation. At genus~$1$, the $S$ and $T$ matrices act on
the $k{+}1$-dimensional space of level-$k$ characters; this $S/T$
action is the level-$5$ modular reconstruction of the open
factorization category, not an intrinsic closed-algebra modularity.
The closed-shadow modular consequences (vector-valued modular
covariance of characters, Verlinde rationality) follow from the
trace + clutching on $\overline{\cM}_{1,1}$ together with named
finiteness/positivity hypotheses (\textsc{C}$_2$-cofiniteness with
Zhu rationality~\cite{Zhu96}; or HS-sewing of
Definition~\ref{def:hs-sewing}); on the \emph{KSDual} sublocus the
reconstruction is an equivalence rather than an adjunction with a
comparison map.

The \emph{first Chern class} of $\mathcal{V}_{g,0}$
(via the Quillen determinant) recovers
$\kappa(\cA) \cdot \lambda_g$ on the diagonal scalar projection, while the
\emph{full Chern character}
$\mathrm{ch}(\mathcal{V}_{g,0})$ recovers the spectral
discriminant~$\Delta_{\cA}$
(Remark~\ref{rem:spectral-characteristic}).
Computing the flat connection (beyond its characteristic
classes) constitutes the passage from Level~$0$ to Level~$2$ of the
modular invariant hierarchy
(Remark~\ref{rem:four-levels},
Conjecture~\ref{conj:categorical-modular-kd}).
\end{remark}

\begin{spectralsequence}[Genus spectral sequence]\label{ss:genus}
The genus spectral sequence has the form:
\[E_2^{p,q} = H^p(\overline{\mathcal{M}}_g, \mathcal{H}^q(\bar{B}^{(g)}(\mathcal{A}))) \Rightarrow \mathbb{H}^{p+q}(\overline{\mathcal{M}}_g, \bar{B}^{(g)}_\bullet(\mathcal{A}))\]
where $\mathbb{H}^*$ denotes hypercohomology.
This is the Leray spectral sequence at fixed genus~$g$; the
all-genus object is the genus-adic inverse limit of finite windows in
Theorem~\ref{thm:master-tower}, with $g_s$ a formal genus parameter
and no asserted analytic convergence.

The differentials have explicit descriptions:
\begin{enumerate}
\item $d_2$: Kodaira--Spencer map (infinitesimal deformations)
\[d_2 = \sum_{i=1}^{3g-3} \frac{\partial}{\partial t_i} \otimes \mu_i\]
where $t_i$ are local coordinates on $\mathcal{M}_g$

\item $d_3$: Massey products (higher compositions).
The differential $d_3$ is induced by the $A_\infty$ operation $m_3$ via the homotopy
transfer data; it takes the form of a triple Massey product in the $E_2$ page

\item $d_r$ ($r \geq 4$): Higher obstruction classes. The differential $d_r\colon E_r^{p,q} \to E_r^{p+r,q-r+1}$ detects higher extension data in the sheaf cohomology of $\overline{\mathcal{M}}_g$. For Koszul algebras, the weight filtration forces these to vanish at bounded stage (Theorem~\ref{thm:bar-cobar-spectral-sequence})
\end{enumerate}

At each fixed genus~$g$, this spectral sequence converges under the standard boundedness condition on $\overline{\mathcal{M}}_g$ (which is compact). The total genus expansion $\sum_g g_s^{2g-2} H^*(\bar{B}^{(g)})$ is a formal power series in $g_s$; its convergence is a separate question related to the Borel summability of the perturbative expansion.
\end{spectralsequence}

% ==========================================
% TOPOLOGICAL RECURSION
% ==========================================

\subsection{Topological recursion and computational methods}\label{subsec:recursion}
\index{topological recursion}

\begin{conjecture}[Eynard--Orantin recursion for bar complex; \ClaimStatusConjectured]\label{conj:EO-recursion}
\index{Eynard--Orantin recursion}
The bar complex correlation functions $\omega_{g,n}$ satisfy the recursion:
\[\omega_{g,n}(z_1,\ldots,z_n) = \sum_{p \in \mathrm{Ram}} \operatorname{Res}_{q \to p} K(z_1, q) \Bigl[\omega_{g-1,n+1}(q,\sigma(q),z_2,\ldots,z_n)\]
\[+ {\sum_{\substack{g_1+g_2=g \\ I \sqcup J = \{2,\ldots,n\}}}}' \omega_{g_1,|I|+1}(q,z_I) \cdot \omega_{g_2,|J|+1}(\sigma(q),z_J)\Bigr]\]

where:
\begin{itemize}
\item $K(z_1,q) = \frac{\frac{1}{2}\int_{\sigma(q)}^q B(z_1, \cdot)}{\omega_{0,1}(q) - \omega_{0,1}(\sigma(q))}$ is the recursion kernel, with $B$ the Bergman kernel (bidifferential of the second kind)
\item $\mathrm{Ram}$ denotes the set of ramification points of the spectral curve
\item $\sigma(q)$ is the conjugate point under the local deck involution near $p$ (not complex conjugation)
\item The prime on $\sum'$ excludes unstable terms $(g_i, |I_i|+1) = (0,1)$
\end{itemize}

Specifically:
\begin{enumerate}
\item[\textup{(a)}] \textup{[\ClaimStatusProvedElsewhere]}
\emph{Heisenberg.} For the free boson, $\omega_{g,n}$ are Gaussian
matrix model correlators and the recursion holds
\cite{Eyn04,CEO06,EO07}.

\item[\textup{(b)}] \textup{[\ClaimStatusConditional{}, conditional on
Conjecture~\textup{\ref{conj:v1-bar-worldline}}; separated from the surrounding conjecture]}
\emph{Koszul chiral algebras.} For a Koszul chiral
algebra~$\mathcal{A}$, the bar complex correlation functions
satisfy the EO recursion (see proof below), conditional on the
propagator--bidifferential correspondence
\textup{(}Conjecture~\textup{\ref{conj:v1-bar-worldline}},
\ClaimStatusHeuristic\textup{)}.

\item[\textup{(c)}] \textup{[\ClaimStatusConjectured]}
\emph{General case.} For a non-Koszul chiral algebra, the
recursion requires verification of the loop equations in the
abstract topological recursion framework.
\end{enumerate}

\end{conjecture}

\begin{remark}[Conditional EO verification]
Let $\mathcal{A}$ be a Koszul chiral algebra equipped with the
spectral-curve comparison of
Conjecture~\ref{conj:v1-bar-worldline} and the all-genus graph
expansion of Theorem~\ref{thm:mk-general-structure-vol1}. Under
these two hypotheses the abstract topological recursion framework
(Kontsevich--Soibelman 2010; Andersen--Borot--Orantin 2017) is
verified as follows.

\emph{Spectral curve from genus-$0$ data.}
The genus-$0$ bar complex data defines a spectral curve
$(S, x, y, B)$: the tree-level OPE coefficients determine
$\omega_{0,1}$ and the ramification data of the projection
$x\colon S \to \mathbb{P}^1$; the bidifferential
$\omega_{0,2} = B$ is the Bergman kernel on~$S$, identified with
the genus-$0$ propagator $G_0(z,w) = 1/(z-w)$ via
Conjecture~\ref{conj:v1-bar-worldline}. The Koszul property (quadratic
OPE relations) ensures that $S$ is a branched cover with finitely
many ramification points.

\emph{Feynman transform decomposition.}
By the assumed all-genus $m_k$ graph expansion, the bar complex
correlation functions are computed by the Feynman transform of the
modular operad $\{\overline{\mathcal{M}}_{g,n}\}$:
\[
\omega_{g,n}(z_1, \ldots, z_n) \;=\; \sum_{\Gamma \in \mathcal{G}_{n,g}}
\frac{1}{|\mathrm{Aut}(\Gamma)|}
\prod_{e \in E(\Gamma)} G(z_{s(e)}, z_{t(e)})
\cdot \prod_{v \in V(\Gamma)} C_v
\]
where the sum is over stable graphs of type~$(g,n)$,
$G$ is the Arakelov--Green propagator, and $C_v$ are OPE structure
constants.

\emph{Edge contraction.}
Every stable graph $\Gamma \in \mathcal{G}_{n,g}$ with $n \geq 1$
has at least one edge~$e_1$ incident to the half-edge labelled by~$z_1$.
Summing over the choice of~$e_1$ and contracting gives a bijection
between $\mathcal{G}_{n,g}$ (with $e_1$ marked) and pairs consisting
of either:
\begin{itemize}
\item a \emph{non-separating} contraction, contributing the
 $\omega_{g-1,n+1}$ term (genus drops by~$1$, two new
 half-edges labelled $q$ and $\sigma(q)$), or
\item a \emph{separating} contraction, contributing
 $\omega_{g_1,|I|+1} \cdot \omega_{g_2,|J|+1}$ with
 $g_1+g_2 = g$ and $I \sqcup J = \{2,\ldots,n\}$.
\end{itemize}
The propagator~$G$ contributes $\operatorname{Res}_{q \to p}$
(extracting the polar part at the ramification point~$p$), and
the genus-$0$ vertex at the contracted edge is identified with the
recursion kernel~$K(z_1, q)$ via the spectral curve data above.
This is the EO recursion.

\emph{Axiom verification.}
We verify the three axioms of the abstract framework:
\begin{itemize}
\item \emph{Symplecticity}: $\omega_{0,2} = B$ is the Bergman kernel,
symmetric with a unique double pole on the diagonal
(Definition~\ref{def:bar-differential-complete}).

\item \emph{Dilaton equation}: The stability condition
$2g - 2 + n > 0$ on the Feynman transform
(Definition~\ref{def:feynman-transform}) forces
\[
\textstyle\sum_{p} \operatorname{Res}_{q \to p} \Phi(q)\,
\omega_{g,n+1}(q, z_1, \ldots, z_n)
= (2g{-}2{+}n)\,\omega_{g,n}(z_1,\ldots,z_n)
\]
where $\Phi = \int^z y\,dx$ is the primitive.

\item \emph{Loop equations}: The constraint $\Dg{g}^{\,2} = 0$
(Theorem~\ref{thm:genus-induction-strict},
Convention~\ref{conv:higher-genus-differentials}) restricted to the
genus-$g$ sector yields quadratic relations among the
$\omega_{g',n'}$ for $2g'+n' < 2g+n$. These are precisely
the loop equations: the total corrected differential decomposes
into a ``residue'' part (extracting OPE poles) and a
``degeneration'' part (contracting edges), and $\Dg{g}^{\,2} = 0$ says
the composition of residue-then-degenerate equals
degenerate-then-residue.
\end{itemize}

By the uniqueness theorem for abstract topological recursion
(Eynard--Orantin 2007; Kontsevich--Soibelman 2010), any collection
of multidifferentials satisfying these axioms with given spectral
curve data is uniquely determined by the recursion formula.
The recursion is the scalar shadow of the MC tautological relation (Corollary~\ref{cor:topological-recursion-mc-shadow}).
\end{remark}

\begin{remark}[Scope]
Part~(a) is the classical result: for the Heisenberg algebra, the
bar complex correlation functions are Gaussian matrix model correlators,
and the EO recursion is a theorem (Eynard 2004;
Chekhov--Eynard--Orantin 2006). Part~(b) is conditional for Koszul
chiral algebras: it requires both the all-genus graph expansion
\textup{(}Theorem~\ref{thm:mk-general-structure-vol1},
\ClaimStatusConjectured\textup{)} and the spectral-curve
propagator--bidifferential comparison
\textup{(}Conjecture~\ref{conj:v1-bar-worldline},
\ClaimStatusHeuristic\textup{)}. Under those two inputs, the
remaining steps are the abstract topological-recursion argument.

Part~(c) remains conjectural: for a non-Koszul chiral
algebra, the spectral curve may not be well-defined (the OPE relations
need not be quadratic), and the loop equations require independent
verification. The abstract topological recursion framework
(Kontsevich--Soibelman 2010; Andersen--Borot--Orantin 2017) provides
the axioms, but their verification for non-Koszul algebras remains open.
The full Maurer--Cartan resolution is Theorem~\ref{thm:mc2-full-resolution}.
\end{remark}

\begin{construction}[Higher genus bar differentials]\label{con:higher-genus-bar}
Given a chiral algebra $\mathcal{A}$, genus $g$, and degree $n$, the total corrected bar differential $\Dg{g}$ \textup{(}Convention~\textup{\ref{conv:higher-genus-differentials})} in degree~$n$ and the correlation $\omega_{g,n}$ are computed recursively.

\emph{Initialization.}
Set $\omega_{0,n}$ to be the tree-level OPE coefficients and let $\mathcal{K}$ denote the Bergman kernel on the spectral curve.

\emph{Recursive computation.}
For $h = 1, \ldots, g$: let $\Omega_h$ be the period matrix of $\Sigma_h$ and $\theta[\alpha]$ the theta functions with characteristics $\alpha$. At each ramification point $p$, update
\[\omega_{h,n} \leftarrow \omega_{h,n} + \operatorname{Res}_{z \to p}\!\left[K(z,p) \cdot \omega_{h-1,n+2} + \sum_{h_1+h_2=h} \omega_{h_1,|I|+1} \cdot \omega_{h_2,|J|+1}\right]\]
and set $d^{(h)}_n = d^{(h-1)}_n + \sum_{\gamma \in H_1(\Sigma_h)} \oint_\gamma \omega_{h,n}$.

\emph{Verification.}
The identity $(d^{(g)}_n)^2=0$ is the genus-$g$ component of the
Maurer--Cartan equation for the corrected differential
(Theorem~\ref{thm:genus-induction-strict}); geometrically it is the
boundary-of-boundary cancellation on
$\overline{\mathcal{M}}_{g,n}$, with curvature terms included in
Convention~\ref{conv:higher-genus-differentials}.
\end{construction}

\begin{example}[Explicit calculation: \texorpdfstring{$\beta\gamma$}{beta-gamma} system through genus 3]
For the $\beta\gamma$ system with $\lambda = 1/2$ (conformal weight):

\emph{Genus~0.} Standard OPE
\[\omega_{0,2} = \frac{dz_1 dz_2}{(z_1 - z_2)^2}\]

\emph{Genus~1.} Elliptic propagator
\[\omega_{1,1} = \frac{c_{\beta\gamma}}{24} \cdot \frac{E_2(\tau)}{2\pi i}\, dz, \qquad c_{\beta\gamma} = 2(6\lambda^2 - 6\lambda + 1)\]
where $E_2(\tau) = 1 - 24\sum_{n=1}^\infty \frac{nq^n}{1-q^n}$ is the second Eisenstein series. At $\lambda = 1/2$, the central charge is $c_{\beta\gamma} = -1$ and the coefficient is $-1/24$.

\emph{Genus~2.} Let $\Omega \in \mathfrak{H}_2$ be the period matrix of~$\Sigma_2$. At $\lambda = 1/2$, the $\beta\gamma$ partition function is:
% The product below runs over the 10 even half-characteristics only;
% the 6 odd characteristics have vanishing theta nulls and are excluded.
\[Z_2(\beta\gamma) = \frac{1}{\det(\mathrm{Im}\,\Omega)^{1/2}} \cdot \frac{1}{\prod_{\delta\,\text{even}}\, \vartheta[\delta](0|\Omega)}\]
where the product runs over the $10$ \emph{even} half-characteristics $\delta \in (\tfrac{1}{2}\mathbb{Z}/\mathbb{Z})^4$ (genus-$2$ theta functions with characteristics; the $6$ odd characteristics have $\vartheta[\delta](0|\Omega) = 0$ and do not appear). The genus-$2$ bar complex differential involves the Siegel--Eisenstein series $E_2(\Omega)$ of degree~$2$ (a function on~$\mathfrak{H}_2$, not to be confused with the genus-$1$ Eisenstein series).

\emph{Genus~3.} The Fay trisecant identity on the Jacobian $\mathrm{Jac}(\Sigma_3) = \mathbb{C}^3/(\mathbb{Z}^3 + \Omega\mathbb{Z}^3)$ provides relations among the genus-$3$ bar complex generators. The Schottky relation (the vanishing of a specific Siegel modular form at the Jacobian locus $\mathcal{J}_3 \subset \mathcal{A}_3$) yields a nontrivial constraint on the $\beta\gamma$ genus-$3$ obstruction, reducing the number of independent tautological integrals.

The pattern is that each genus~$g$ introduces degree-$g$ Siegel modular forms. For the bosonic $\beta\gamma$ system, the partition function depends on the choice of line bundle $K_X^\lambda$, not on spin structures; spin structure dependence arises instead for fermionic systems such as the $bc$ ghost.
\end{example}

% ==========================================
% PHYSICAL INTERPRETATION
% ==========================================

\subsection{Physical interpretation: strings and holography}\label{subsec:physics}

\begin{theorem}[Bar complex defines tautological moduli integrals; \ClaimStatusProvedHere]\label{thm:bar-moduli-integrals}
\textbf{Type signature.} (Open quadrant; level~$5$ scalar shadow on
$\overline{\cM}_{g,n}$ with module insertions; hypothesis package:
Koszul chiral algebra $\cA$ with modules
$\mathcal{V}_1,\ldots,\mathcal{V}_n$; Feynman-transform identification
\textup{(}Theorem~\textup{\ref{thm:prism-higher-genus}}\textup{)};
properness of the Fulton--MacPherson configuration map.)
For a Koszul chiral algebra $\mathcal{A}$ and modules
$\mathcal{V}_1,\ldots,\mathcal{V}_n$, the bar complex with module
insertions defines a tautological class
\[
\operatorname{cl}_{g,n}(\mathcal{V}_1,\ldots,\mathcal{V}_n)
\in H^\bullet(\overline{\mathcal{M}}_{g,n})
\]
by proper pushforward from the logarithmic Fulton--MacPherson
configuration space.  Pairing this class with the virtual fundamental
class gives
\[
\mathcal{I}_{g,n}(\mathcal{V}_1,\ldots,\mathcal{V}_n)
=
\int_{[\overline{\mathcal{M}}_{g,n}]}
\operatorname{cl}_{g,n}(\mathcal{V}_1,\ldots,\mathcal{V}_n).
\]
The genus expansion
\[\mathcal{I}_{\mathrm{total}} = \sum_{g=0}^\infty g_s^{2g-2+n}\, \mathcal{I}_{g,n}(\mathcal{V}_1,\ldots,\mathcal{V}_n)\]
is the formal power series in $g_s$ assembled from the genus-graded bar complex (Theorem~\ref{thm:master-tower}).
\end{theorem}

\begin{proof}
The bar complex at genus~$g$ with module insertions is a logarithmic
sheaf on the compactified configuration space over
$\overline{\mathcal{M}}_{g,n}$, constructed via the Feynman transform
identification
$\bar{B}^{(g)}_n(\mathcal{V}_1 \otimes \cdots \otimes \mathcal{V}_n)
\simeq
\mathrm{FT}_{\mathcal{M}\mathrm{od}}(\mathcal{A};
\mathcal{V}_1,\ldots,\mathcal{V}_n)_g$
(Theorem~\ref{thm:prism-higher-genus}).  The forgetful map
$\pi\colon \overline{C}_n^{(g)}(\Sigma_g)\to
\overline{\mathcal{M}}_{g,n}$ is proper
(Definition~\ref{def:higher-genus-config}), so
$R\pi_*$ sends the bar cohomology class to a class on
$\overline{\mathcal{M}}_{g,n}$.  Since
$\overline{\mathcal{M}}_{g,n}$ is a proper Deligne--Mumford stack,
pairing with its virtual fundamental class is defined.  The formal
genus assembly follows from the finite-window genus tower of
Theorem~\ref{thm:master-tower}.
\end{proof}

\begin{conjecture}[String amplitude = bar complex cohomology; \ClaimStatusConjectured]\label{conj:string-amplitude-bar}
For a chiral algebra $\mathcal{A}$ describing string theory vertex operators, the moduli integrals of Theorem~\ref{thm:bar-moduli-integrals} compute string scattering amplitudes:
\[\mathcal{A}_{g,n}^{\mathrm{string}}(\mathcal{V}_1,\ldots,\mathcal{V}_n) = \mathcal{I}_{g,n}(\mathcal{V}_1,\ldots,\mathcal{V}_n)\]
where $\mathcal{A}_{g,n}^{\mathrm{string}}$ is the $g$-loop, $n$-point string amplitude and $\mathcal{V}_i$ are string vertex operators. The string coupling expansion $\mathcal{A}^{\mathrm{string}}_{\mathrm{total}} = \sum_{g=0}^\infty g_s^{2g-2+n}\, \mathcal{A}_{g,n}^{\mathrm{string}}$ matches the genus expansion of the bar complex.
\end{conjecture}

\begin{remark}[Scope]
The moduli integrals $\mathcal{I}_{g,n}$ are established in Part~\ref{part:bar-complex}.
At genus~$0$, the algebraic BRST/bar comparison is proved, while the
moduli-integration/amplitude identification is only conditional
\textup{(}Corollary~\ref{cor:string-amplitude-genus0},
Theorem~\ref{thm:brst-bar-genus0}\textup{)}; for affine
Kac--Moody algebras, $\barB(\cA) \simeq C^{\infty/2+*}(\cA)$ at all
levels $k \neq -h^\vee$
\textup{(}Theorem~\ref{thm:bar-semi-infinite-km}\textup{)}; and the
principal W-algebra extension is conditional on
Theorem~\ref{thm:bar-semi-infinite-w}. Extension to all genera
requires the higher-genus BRST-bar identification
\textup{(}Conjecture~\ref{conj:anomaly-physical}\textup{)}; at
genus~$0$ the algebraic BRST-bar comparison is proved, while the
full genuswise BV/BRST/bar identification remains open
at genus~$\geq 1$
\textup{(}Conjecture~\ref{conj:master-bv-brst}\textup{)}. At genus~$\geq 2$, the algebraic bar differential is determined
inductively (Theorem~\ref{thm:inductive-genus-determination});
the analytic graph amplitudes converge without UV renormalization
in dimension~$2$ (Proposition~\ref{prop:2d-convergence}); and the
analytic-algebraic comparison holds on the algebraic core
(Theorem~\ref{thm:analytic-algebraic-comparison}).
The remaining input for the full analytic identification is the
HS-sewing condition (Definition~\ref{def:hs-sewing}) for mode-sum
convergence.
\end{remark}


\begin{conjecture}[Holographic duality via bar construction; \ClaimStatusConjectured]
\label{conj:holographic-bar-cobar}
The pair $(\cA,\barB(\cA))$ supplies an algebraic comparison datum for
AdS/CFT: OPE coefficients correspond to vertices, conformal blocks to
Witten diagrams, $c$ to $R_{\mathrm{AdS}}$, and $1/N$ to genus.  The
cobar inversion $\Omega\barB(\cA)\simeq\cA$ is the boundary
resolution input; the Koszul dual algebra $\cA^!$ and the derived
centre $Z_{\mathrm{ch}}^{\mathrm{der}}(\cA)$ are separate objects.
The conjectural comparison
$Z_{\mathrm{CFT}}[\mathcal{J}]=Z_{\mathrm{gravity}}[\phi_0=\mathcal{J}]$
identifies the genus expansion with the $1/N$ expansion
($g\mapsto N^{-(2g-2)}$). It requires bulk Chern--Simons
reconstruction~\cite{Wit89}, Witten-diagram matching, and the
genuswise BV/BRST/bar comparison of
Conjecture~\ref{conj:master-bv-brst}; the genus-$0$ and genus-$1$
algebraic comparison surfaces are the proved part recorded in the
concordance.
\end{conjecture}

% ==========================================
% MODULAR HIERARCHY
% ==========================================

\subsection{Modular hierarchy}\label{subsec:modular-hierarchy}

\begin{theorem}[Poincaré--Verdier duality for genus pieces; \ClaimStatusProvedElsewhere{} \cite{BD04,FG12}]\label{thm:poincare-extended}
\index{Poincare-Verdier duality@Poincar\'e--Verdier duality!higher genus}
\textbf{Type signature.} (Open quadrant; Verdier-Koszul reflection at
level~$2$, with the curve-Verdier shift on $\overline{\cM}_{g,n}$;
hypothesis package: $\cA^i = H^\bullet B^{\mathrm{ch}}(\cA)$
finite-type; perfect genus-$g$ Poincar\'e--Koszul pairing.)
Let $B^{(g)}(\mathcal{A})$ be the genus-$g$ bar object and let
$\mathcal{A}^i:=H^\bullet B^{\mathrm{ch}}(\mathcal{A})$ be the
Koszul dual coalgebra. If $\mathcal{A}^i$ is finite-type and the
genus-$g$ Poincar\'e--Koszul pairing between the bar coefficients of
$\mathcal A$ and the cobar coefficients of $\mathcal A^i$ is perfect,
then Verdier duality on $\overline{\mathcal{M}}_{g,n}$ identifies the
dual genus piece with the cobar object built from $\mathcal{A}^i$:
\[
\mathbb D_{\overline{\mathcal{M}}_{g,n}}
\bigl(\bar{B}^{(g)}(\mathcal{A})\bigr)
\cong
\Omega^{(g)}(\mathcal{A}^i).
\]
This is Verdier-dual Koszul passage. It is not the inversion
$\Omega(B(\mathcal{A}))\simeq\mathcal{A}$ and it is not a statement
about the derived centre $Z^{\mathrm{der}}_{\mathrm{ch}}(\mathcal{A})$.
\end{theorem}

\begin{proof}
The bar complex at genus $g$ uses logarithmic forms on
$\overline{\mathcal{M}}_{g,n}$ with coefficients in
$j_*j^*\mathcal{A}^{\boxtimes n}$
(Definition~\ref{def:geometric-bar}). Verdier duality exchanges the
logarithmic extension with the dual logarithmic extension and reverses
the residue maps. The perfect Poincar\'e--Koszul pairing identifies
the Verdier-dual coefficient system of the $\mathcal A$-bar object
with the cobar coefficient system generated by $\mathcal A^i$; finite
type gives $\mathcal A^!=(\mathcal A^i)^\vee$ on the algebra side.
The reversed residue differential is exactly the cobar differential.
Thus the Verdier dual of the genus-$g$ bar object is the genus-$g$
cobar object on $\mathcal{A}^i$. The dualizing-complex shift is
contained in $\omega^\bullet_{\overline{\mathcal{M}}_{g,n}}$. The
construction is compatible with the Feynman transform of the modular
operad
$\{\overline{\mathcal{M}}_{g,n}\}$
(Theorem~\ref{thm:prism-higher-genus}).
\end{proof}

% ==========================================
% THE HOLOGRAPHIC DICTIONARY AND MODULAR STATE SPACES
% ==========================================

\subsection{The holographic dictionary and modular state spaces}
\label{subsec:holographic-dictionary}
\index{holographic dictionary!chain-level}

The chain-level modular functor of the preceding subsection has an
algebraic open--closed reading: the universal acting bulk is the
chiral derived centre of the boundary
$Z^{\mathrm{der}}_{\mathrm{ch}}(\cA_\partial)$ (level~$3$ in the
Beilinson tower), while identification with a physical
three-dimensional holomorphic-topological bulk is a separate
comparison datum (the level-$3$ vertical equivalence to
$\mathsf{SC}^{\mathrm{ch,top}}$-brace bulk in Vol~II). Factorization
homology then produces genus-graded state spaces equipped with a
formal partition function whose $\mathrm{SL}(2,\Z)$-equivariance is
the level-$5$ modular reconstruction (trace + clutching on the open
factorization category), not an intrinsic closure datum.
The complete packaging of these ingredients (bulk algebra, Koszul-dual
defect algebra, algebraic line-side model, genus-zero $r$-matrix,
genus-refined MC class, and shadow connection) is the
\emph{holographic modular Koszul datum} $\mathcal{H}(T)$
(Definition~\ref{def:holographic-modular-koszul-datum}). The modular
$r$-matrix of Definition~\ref{def:modular-r-matrix-genus} below is the
genus-refined specialization of the binary collision residue
$\operatorname{Res}^{\mathrm{coll}}_{0,2}(\Theta_\cA) = r(z)$
(equation~\eqref{eq:collision-residue-identification}). The Vol~III
two-stage cross-volume bridge is
$\Phi_d^{(\Sigma_{d-1},C)}=\mathrm{Sp}^{\mathrm{ch}}_{\Sigma_{d-1},C}\circ
\Phi_d^{\mathrm{FA}}$ at level~$2$; its restriction to
$C=E_\tau$ recovers the elliptic specialization used in
Theorem~\ref{thm:elliptic-bar} above.

\subsubsection{The holographic/Koszul dictionary}

\begin{proposition}[Algebraic closed sector from the boundary;
\ClaimStatusProvedElsewhere{} \cite{BD04,FG12,CG17}]
\label{prop:bulk-from-boundary}
\index{bulk-boundary correspondence!derived center}
\textbf{Type signature.} (Open quadrant; level~$1\to 3$ reconstruction
$A_b \mapsto Z^{\mathrm{der}}_{\mathrm{ch}}(A_b)$ via chiral
Hochschild cochains; hypothesis package: factorisation dg-category
$\mathcal{C}^{\mathrm{op}}$ on $(X,D,\tau)$ with vacuum $b$;
$\cA_\partial = \End_{\mathcal{C}}(b)$; Hochschild cochain complex
on the chiral category in the sense of BD/FG/CG.)
Let $\cA_\partial$ be a boundary chiral algebra. Its universal
algebraic closed-sector object is the chiral derived centre. Thus
\begin{equation}\label{eq:bulk-derived-center}
 \cA_{\mathrm{cl}}^{\mathrm{alg}}
 \;\simeq\;
 Z_{\mathrm{ch}}^{\mathrm{der}}(\cA_\partial),
\end{equation}
where $Z_{\mathrm{ch}}^{\mathrm{der}}$ denotes chiral Hochschild
cochains of the chiral category. Identification with a physical
three-dimensional holomorphic-topological bulk is an additional
comparison datum, not part of this proposition.
\end{proposition}

\begin{remark}[Maurer--Cartan moduli and defect couplings]
\label{rem:mc-defect-couplings}
\index{Maurer--Cartan!defect couplings}
Given a Koszul dual pair $(\cA, \cA^!)$, the Maurer--Cartan moduli
of the completed tensor product classifies consistent defect couplings:
\begin{equation}\label{eq:mc-defect}
 \mathrm{MC}\bigl(\cA^! \,\widehat{\otimes}\, R\bigr)
 \;\longleftrightarrow\;
 \operatorname{Hom}(\cA_\partial, B),
\end{equation}
where $R$ is the formal neighborhood of a classical solution and $B$
is the target boundary condition. This is the chain-level
implementation of defect coupling: each Maurer--Cartan element
specifies a coherent coupling of closed-sector fields to boundary data.
\end{remark}

\subsubsection{Modular state spaces via factorization homology}

\begin{definition}[Genus-graded holographic state spaces]
\label{def:holographic-state-spaces}
\index{factorization homology!holographic state spaces}
Let $L_1, \dotsc, L_n \in \mathcal{C}_T$ be objects of the chosen
line-side model for the $3$d theory $T$ with boundary condition $B$.
The
\emph{holographic state space} at genus $g$ with $n$ insertions is
\begin{equation}\label{eq:holographic-state-space}
 \mathcal{H}_{g,n}^{T;\,B}(L_1, \dotsc, L_n)
 \;:=\;
 \int_{\Sigma_{g,n}}
 \mathrm{Fact}_{T;\,B}[L_1, \dotsc, L_n],
\end{equation}
where the integral denotes factorization homology over the
genus-$g$ surface $\Sigma_{g,n}$ with $n$ boundary punctures.
\end{definition}

\begin{definition}[Genus-graded holographic partition function]
\label{def:holographic-partition-function}
\index{partition function!genus-graded holographic}
The all-genera holographic partition function is
\begin{equation}\label{eq:holographic-pf}
 Z_{T;\,B}(\hbar, t)
 \;:=\;
 \sum_{g,n \ge 0}
 \hbar^{2g-2+n}\;
 \frac{1}{n!}\;
 \chi_{\mathrm{gr}}\bigl(
 \mathcal{H}_{g,n}^{T;\,B}(t)\bigr),
\end{equation}
where $\chi_{\mathrm{gr}}$ is the graded Euler characteristic and
$t$ is a formal insertion parameter. This is the algebraic
all-genera avatar of the large-$N$ boundary/bulk comparison, with
$\hbar$ tracking the genus and $t$ tracking the number of insertions.
\end{definition}

\subsubsection{The modular Yangian kernel}

\begin{definition}[Modular $r$-matrix]
\label{def:modular-r-matrix-genus}
\index{r-matrix@$r$-matrix!modular genus-refined}
The \emph{modular Yangian kernel} is the genus-refined $r$-matrix
\begin{equation}\label{eq:modular-r-genus}
 r_T^{\mathrm{mod}}(z, \hbar)
 \;=\;
 \sum_{g \ge 0} \hbar^{2g}\, r_{T,g}(z),
 \qquad
 r_{T,0}(z) = r_T(z),
\end{equation}
where each coefficient lives in
$H^\bullet(\overline{\mathcal{M}}_{g,2}) \otimes
\cA_T^! \,\widehat{\otimes}\, \cA_T^!\bigl(\!(z^{-1})\!\bigr)$.
At $g = 0$ this is the dg-shifted Yangian $r$-matrix of
Definition~\ref{def:dg-shifted-yangian}; the modular completion is a
genus-refined, curve-valued OPE kernel.
\end{definition}

\subsubsection{Descendant holographic amplitudes}

\begin{definition}[Descendant holographic amplitudes]
\label{def:descendant-hol-amplitudes}
\index{holographic amplitudes!descendant}
For operators $O_1, \dotsc, O_n$ and non-negative integers
$k_1, \dotsc, k_n$, the \emph{descendant holographic amplitude} at
genus $g$ is
\begin{equation}\label{eq:descendant-amplitude}
 \bigl\langle
 \tau_{k_1}(O_1) \dotsm \tau_{k_n}(O_n)
 \bigr\rangle_g^{\mathrm{hol}}
 \;:=\;
 \int_{[\mathrm{FM}_n(\Sigma_g \,|\, D)]}
 \Bigl(\prod_{i=1}^n \psi_i^{k_i}\Bigr)\;
 \omega_{O_1, \dotsc, O_n},
\end{equation}
where $\mathrm{FM}_n(X \,|\, D)$ denotes the logarithmic
Fulton--MacPherson space compactifying configuration space relative to
the divisor $D$, and $\psi_i$ are the cotangent-line classes. The
descendant amplitudes connect the chain-level modular functor to the
Gromov--Witten-type intersection theory of
\S\ref{sec:modular-chern-weil-transform}.
\end{definition}

\begin{remark}[Finite data]\label{rem:finite-data-genus}
\index{genus expansion!finite data}
The diagonal scalar genus expansion is determined by the genus-$1$
curvature coefficient:
\[
F_g^{\mathrm{diag}}(\mathcal{A})
=\kappa(\mathcal{A})\lambda_g^{\mathrm{FP}}.
\]
On the uniform-weight scalar lane
(Definition~\ref{def:scalar-lane}) this diagonal term is the full
free energy at all genera. For arbitrary modular Koszul algebras the
diagonal term is still defined at every genus, but the full
all-weight free energy is
$F_g=F_g^{\mathrm{diag}}+\delta F_g^{\mathrm{cross}}$, with
$\delta F_1^{\mathrm{cross}}=0$ universally and
$\delta F_g^{\mathrm{cross}}$ generally nonzero at $g\geq2$.

The standard scalar coefficients are
\[
\kappa(V_k(\mathfrak{g}))
=\frac{(k+h^\vee)\dim\mathfrak{g}}{2h^\vee},
\qquad
\kappa(\mathrm{Vir}_c)=\frac{c}{2},
\qquad
\kappa(\mathcal{H}_k)=k,
\]
and, for principal $\mathcal{W}_N$,
$\kappa(\mathcal{W}_N)=c(H_N-1)$.  For $\mathcal{W}_3$ the first
mixed correction is
$\delta F_2^{\mathrm{cross}}=(c+204)/(16c)$
(Theorem~\ref{thm:multi-weight-genus-expansion}). The spectral
invariants and the full Maurer--Cartan completion are not determined
by~$\kappa(\cA)$ alone; see Theorem~\ref{thm:spectral-characteristic}
and Definition~\ref{def:full-modular-package}.
The Mumford $\lambda$-class relations are projections of the MC equation at degree~$2$ (Proposition~\ref{prop:mumford-from-mc}).
\end{remark}


%================================================================
% THE MODULAR CHERN--WEIL TRANSFORM AND GEOMETRIC AVATARS
%================================================================

\section{The modular Chern--Weil transform}
\label{sec:modular-chern-weil-transform}

The diagonal scalar expansion
$F_g^{\mathrm{diag}}(\cA)=\kappa(\cA)\lambda_g^{\mathrm{FP}}$, the
mixed-channel correction $\delta F_g^{\mathrm{cross}}$, and the
spectral discriminant $\Delta_{\cA}$ are characteristic images of one
Maurer--Cartan element. The modular Chern--Weil transform is a single
map from the ambient master element to the full invariant hierarchy.

\begin{definition}[Modular Chern--Weil transform]
\label{def:modular-chern-weil-transform}
Define the \emph{modular Chern--Weil transform}
\begin{equation}
\label{eq:modular-chern-weil-transform}
\operatorname{CW}_{\cA}\colon
\operatorname{MC}(\mathfrak{g}^{\mathrm{amb}}_{\cA})
\longrightarrow
\underbrace{\vphantom{\bigl(}
\text{scalar classes }}_{F_g^{\mathrm{diag}},\ \delta F_g^{\mathrm{cross}}}
\;\times\;
\underbrace{\vphantom{\bigl(}
\text{spectral classes}}_{\Delta_{\cA}(x)}
\;\times\;
\underbrace{\vphantom{\bigl(}
\text{resonance classes}}_{\mathfrak{R}^{\mathrm{mod}}_{r,\bullet}}
\;\times\;
\underbrace{\vphantom{\bigl(}
\text{line-kernel data}}_{r_{\cA}(z)}
\end{equation}
by the rules
\begin{align}
\operatorname{tr}_{\mathrm{diag}}(\Theta_{\cA})
&= \sum_{g \geq 1}\kappa(\cA)\,\lambda_g,
\qquad \label{eq:cw-scalar}
\\
\operatorname{tr}_{\mathrm{cross}}(\Theta_{\cA})
&= \sum_{g \geq 2}\delta F_g^{\mathrm{cross}}(\cA),
\qquad \label{eq:cw-cross}
\\
\det(1 - xT_{\mathrm{br},\cA})
&= \Delta_{\cA}(x),
\label{eq:cw-spectral}
\\
\operatorname{CW}_{\cA}^{\mathrm{res},r}(\Theta_{\cA})
&= \mathfrak{R}^{\mathrm{mod}}_{r,\bullet}(\cA),
\qquad r \geq 4,
\label{eq:cw-resonance}
\\
\operatorname{CW}_{\cA}^{\mathrm{line}}(\Theta_{\cA})
&= r_{\cA}(z)
\in \cA^!\widehat\otimes \cA^!\bigl(\!(z^{-1})\!\bigr).
\label{eq:cw-line}
\end{align}
\end{definition}

\begin{remark}[Analogy with Chern--Weil]
The scalar pair
$(F_\bullet^{\mathrm{diag}},\delta F_\bullet^{\mathrm{cross}})$,
the spectral $\Delta_\cA$, the nonlinear shadow tensors, the
resonance tower, and the line-operator kernel are successive characteristic images of
$\Theta_{\cA}$ under projection to
different invariant-theoretic targets, the modular analogue
of the classical Chern--Weil homomorphism, where a single connection
produces all characteristic classes.
The shadow tautological classes $\tau_{g,n}(\cA)$ (Definition~\ref{def:shadow-tautological-classes}) form a cohomological field theory under this transform (Theorem~\ref{thm:shadow-cohft}; conditional on the flat-identity hypothesis), with the complementarity propagator serving as the Givental $R$-matrix (Theorem~\ref{thm:cohft-reconstruction}).
\end{remark}

\begin{remark}[K-theoretic hierarchy of the Chern--Weil transform]
\label{rem:k-theoretic-hierarchy}
\index{Chern--Weil theory!K-theoretic hierarchy}
\index{modular tangent complex!K-theoretic interpretation}
The Chern--Weil projections
(eq.~\eqref{eq:cw-scalar}--\eqref{eq:cw-line}) correspond
to characteristic invariants of the modular tangent complex
$T^{\mathrm{mod}}_{\Theta_\cA}$
(Construction~\ref{const:vol1-modular-tangent-complex})
at successive levels of nonlinearity:
\begin{center}
\renewcommand{\arraystretch}{1.2}
\begin{tabular}{lll}
\textbf{Characteristic invariant} & \textbf{CW projection}
 & \textbf{Shadow} \\
\hline
First Chern class $c_1$ & Diagonal scalar trace & $\kappa(\cA)$ \\
Mixed boundary character & Cross trace & $\delta F_g^{\mathrm{cross}}$ \\
Chern character & Spectral determinant & $\Delta_\cA(x)$ \\
Higher characteristic classes ($r \geq 4$)
 & Resonance classes & $\mathfrak{R}^{\mathrm{mod}}_{r,\bullet}$ \\
Holonomy representation & Line kernel & $r_\cA(z)$
\end{tabular}
\end{center}
The diagonal scalar projection
$\kappa(\cA) = c_1(\det T^{\mathrm{mod}})$ is the first Chern
class of the determinant line bundle. The cross trace is zero at
genus~$1$, vanishes on the scalar lane, and is generally nonzero for
interacting multi-weight algebras at $g\geq2$. The spectral
projection $\Delta_\cA =
\operatorname{sdet}(1 - t \cdot H(d_{\Theta_\cA}|_{\mathrm{red}}))$
is the Chern character
(eq.~\eqref{eq:delta-chern-character}). The resonance
classes at degree $r \geq 4$ are the first nonlinear
characteristic classes: they carry information not determined
by $\kappa(\cA)$ and $\Delta_\cA$ alone, and are the shadow-tower data
that distinguishes algebras at the same scalar curvature.
\end{remark}

The transform on the two extreme examples is computed next.

\begin{example}[Chern--Weil transform: Heisenberg]
\label{ex:chern-weil-heisenberg}
For $\mathcal{H}_\kappa$ with $\Theta_{\mathcal{H}} = \kappa
\cdot \eta \otimes \Lambda$ (the scalar element):
\begin{enumerate}[label=\textup{(\roman*)}]
\item \emph{Scalar projection}:
$\operatorname{tr}(\Theta_{\mathcal{H}})
= \sum_{g \geq 1} \kappa \cdot \lambda_g^{\mathrm{FP}}$. The genus tower is determined by $\kappa$ alone.
\item \emph{Spectral projection}: The branching operator is trivial
on a rank-$1$ slice: $\Delta_{\mathcal{H}}(x) = 1 - \kappa x$.
One spectral branch at $x = 1/\kappa$.
\item \emph{Resonance tower}: $\mathfrak{R}^{\mathrm{mod}}_r
(\mathcal{H}) = 0$ for all $r \geq 4$. No nonlinear shadows
means no resonance: the contact quartic sector is empty.
\item \emph{Line kernel}: $r_{\mathcal{H}}(z)=\kappa/z$ in rank~$1$
(and $k\,\mathrm{id}_V/z$ in rank~$\dim V$). It is abelian and
rank-one, but not zero unless the level vanishes.
\end{enumerate}
The Heisenberg is the unique standard family where every
Chern--Weil projection is a function of the single scalar
$\kappa$ alone:
$F_g = \kappa\lambda_g^{\mathrm{FP}}$,
$\delta F_g^{\mathrm{cross}}=0$,
$\Delta(x) = 1 - \kappa x$,
$\mathfrak{R}_r = 0$ for all $r \geq 4$, and
$r_{\mathcal{H}}(z)=\kappa/z$. No shadow jets exist
($\mathfrak{C} = \mathfrak{Q} = 0$), so the Chern--Weil
map is rank-$1$ in every channel. By contrast, Virasoro
carries independent data in all four channels.
\end{example}

\begin{example}[Chern--Weil transform: Virasoro]
\label{ex:chern-weil-virasoro}
For $\operatorname{Vir}_c$ with $H = (c/2)x^2$,
$\mathfrak{C} = 2x^3$,
$\mathfrak{Q}^{\mathrm{contact}} = 10/[c(5c{+}22)] \cdot x^4$:
\begin{enumerate}[label=\textup{(\roman*)}]
\item \emph{Scalar}: $\operatorname{tr}(\Theta_{\mathrm{Vir}})
= \sum_g (c/2)\lambda_g^{\mathrm{FP}}$. Same $\hat{A}$-genus as any algebra with $\kappa(\cA) = c/2$.
\item \emph{Spectral}:
$\Delta_{\mathrm{Vir}}(x) = 1 - (c/2)x$ on the primary line.
\item \emph{Resonance}: The quartic resonance class has singular
locus $\{c(5c{+}22) = 0\}$, consisting of:
\begin{itemize}
\item $c = 0$: the uncurved point where $\kappa(\mathrm{Vir}_c) = 0$ and the
Hessian degenerates;
\item $c = -22/5$: the Lee--Yang value $c_{5,2}$, where the
Kac determinant at level $4$ vanishes and $\Lambda_{\mathrm{Vir}}$
becomes null.
\end{itemize}
At these loci, the quartic contact coefficient
$Q^{\mathrm{contact}} = 10/[c(5c{+}22)]$ has a pole:
the shadow geometry degenerates.
\item \emph{Line kernel}: $r_{\mathrm{Vir}}(z)$ is the line-OPE
kernel for Virasoro representations, computed from the
dg-shifted Yangian of $\mathfrak{sl}_2$ at the level
determined by~$c$.
\end{enumerate}
The Virasoro is the simplest family where the scalar, spectral,
resonance, and line channels carry nontrivial data. The scalar and
spectral are identical to
any algebra with $\kappa(\cA) = c/2$, but the resonance channel
distinguishes the Virasoro from other algebras at the same central
charge: $\operatorname{CW}_{\cA}$ carries strictly
more information than $\kappa(\cA)$ alone.
\end{example}

\begin{remark}[Completion entropy is genus-independent]
\label{rem:entropy-genus-independent}
\index{completion entropy!genus-independence}
The completion entropy $h_K(\cA) = \log(\rho_\cA^{-1})$
(Definition~\ref{def:completion-entropy}) is a purely kinematic
invariant: it depends only on the vacuum character $\chi_\cA(q)$,
not on the genus, the curvature $\kappa(\cA)$, or any OPE
structure constants. In the language of the Chern--Weil transform,
$h_K(\cA)$ is a \emph{zero-genus, zero-order} invariant: it
measures the combinatorial growth of the completed bar complex
before any genus correction or shadow data is imposed. The
genus-$g$ corrections (Theorem~D, the shadow obstruction tower,
the modular QME) act on a fixed kinematic substrate whose
complexity is $h_K$.

In particular, the Virasoro involution $c \mapsto 26 - c$ preserves
$h_K(\mathrm{Vir})$ but changes the scalar shadow
$\kappa(\mathrm{Vir}_c) \mapsto 13 - \kappa(\mathrm{Vir}_c)$. The entropy ladder
$h_K(\mathcal{W}_2) < h_K(\mathcal{W}_3) < \cdots <
h_K(\mathcal{W}_\infty)$
(Proposition~\ref{prop:wn-entropy-ladder}) is a \emph{static}
landscape invariant: it does not flow with genus. The Chern--Weil
channels (scalar, spectral, resonance, line) all operate on this
fixed kinematic base.
\end{remark}

\subsection{The complementarity Legendrian}
\label{subsec:complementarity-legendrian}

\begin{definition}[Complementarity Legendrian]
\label{def:complementarity-legendrian}
Contactize the Hamiltonian $\mathfrak{S}_{\cA}$
(Definition~\ref{def:complementarity-action}) to obtain the
Legendrian subspace
\begin{equation}
\label{eq:complementarity-legendrian}
\Lambda_{\cA}
:= j^1\mathfrak{S}_{\cA}
\subset J^1(V_{\cA}).
\end{equation}
The front projection $\pi_{\mathrm{front}}(\Lambda_{\cA})$ is the
graph of $\mathfrak{S}_{\cA}$ in the base coordinates, and its
singularities correspond to degeneracies of the shadow jets.
\end{definition}

\begin{conjecture}[Front-discriminant realization of resonance]
\label{conj:front-discriminant-resonance}
\ClaimStatusConjectured
For each degree $r \geq 4$, the resonance divisor
$\mathfrak{R}^{\mathrm{mod}}_{r,g,n}(\cA)$ is the
front-singularity discriminant of the degree-$r$ contact sector of
$\Lambda_{\cA}$.
\end{conjecture}

\subsection{The complementarity Stokes groupoid and phase index}
\label{subsec:stokes-groupoid-phase-index}

\begin{definition}[Nonlinear regular locus]
\label{def:nonlinear-regular-locus}
Let $B^{\circ}_{\cA}$ be the quartic-regular locus of parameters.
Define
\begin{equation}
\label{eq:nonlinear-regular-locus}
U_{\cA}
:=
B^{\circ}_{\cA} \setminus
\Bigl(
\Delta_{\cA}^{-1}(0)
\cup
\bigcup_{r \geq 4}
|\mathfrak{R}^{\mathrm{mod}}_r(\cA)|
\Bigr).
\end{equation}
\end{definition}

\begin{definition}[Complementarity Stokes groupoid]
\label{def:complementarity-stokes-groupoid}
Let $\mathscr{P}_{\cA}$ denote the local system of gauge classes
of truncated lifts over $U_{\cA}$. Its monodromy groupoid is the
\emph{complementarity Stokes groupoid}
\begin{equation}
\label{eq:complementarity-stokes-groupoid}
\operatorname{Sto}_{\mathrm{comp}}(\cA)
:=
\operatorname{Mon}(\mathscr{P}_{\cA}|_{U_{\cA}}).
\end{equation}
\end{definition}

\begin{definition}[Phase index]
\label{def:phase-index}
For a point $u \in U_{\cA}$, the \emph{nonlinear phase index} is
\begin{equation}
\label{eq:phase-index}
\nu_{\mathrm{nl}}(\cA, u)
:=
\min\{r \geq 3 : \operatorname{Sh}_r(\Theta_{\cA}, u) \neq 0\},
\end{equation}
with $\nu_{\mathrm{nl}} = \infty$ if all higher jets vanish.
\end{definition}

\begin{remark}[Phase index of the standard families]
\label{rem:phase-index-families}
The phase index gives a precise hierarchy:
\[
\begin{array}{c|ccccc}
\text{Family}
& \text{Heisenberg}
& \widehat{\mathfrak{sl}}_2
& \beta\gamma
& \text{Virasoro}
& \text{principal }\mathcal{W}_N
\\[2pt]\hline\rule{0pt}{11pt}
\nu_{\mathrm{nl}}
& \infty
& 3
& 4
& \text{mixed}
& \text{mixed}
\end{array}
\]
The three primitive values $\{3, 4, \infty\}$ correspond to the three
primitive archetypes (Lie/tree, contact/quartic, Gaussian). The
``mixed'' entries for Virasoro and $\mathcal{W}_N$ mean that both
cubic and quartic jets are simultaneously nonzero: these families
have infinite shadow depth, and their
complementarity potentials are non-polynomial
(Theorem~\ref{thm:nms-finite-termination}).
\end{remark}

\begin{remark}[Stokes data and the CohFT]
\label{rem:stokes-r-matrix}
\index{Stokes groupoid!CohFT structure}
The complementarity Stokes groupoid
$\operatorname{Sto}_{\mathrm{comp}}(\cA)$ and the shadow CohFT
$\Omega_{g,n}^{\cA}$ (Theorem~\ref{thm:shadow-cohft};
conditional on the flat-identity hypothesis) are
governed by the same data (the complementarity
propagator~$P_\cA$ and its $\psi$-class expansion) but in
different guises. The Stokes groupoid detects the
\emph{analytic continuation} of the shadow obstruction tower across
the parameter space (with singularities at
$\kappa(\cA) = 0$ and at Kac-determinant zeros), while
the CohFT detects the \emph{tautological projection} of the
MC element to moduli of curves. For class-$\mathbf{G}$
algebras, both are trivial. For class-$\mathbf{M}$, the
non-semisimple locus of the genus-$0$ shadow CohFT
(where the Givental classification fails) coincides with
the singular locus of the complementarity Legendrian
$\Lambda_\cA$ (where the Stokes groupoid has
non-trivial monodromy). A precise identification of
Stokes multipliers with CohFT wall-crossing data
remains open.
\end{remark}

\subsection{The nonlinear modular characteristic hierarchy}
\label{subsec:nonlinear-modular-hierarchy}

\begin{theorem}[The full modular invariant hierarchy]
\label{thm:full-modular-invariant-hierarchy}
\ClaimStatusConditional
\textbf{Type signature.} (Open quadrant; bar-intrinsic MC presentation
on the modular operad $\{\overline{\cM}_{g,n}\}$; spans levels~$1$
\textup{(}chart curvature\textup{)}, $3$ \textup{(}derived chiral
centre / shadow obstruction tower\textup{)}, $4$
\textup{(}line/brane operators via $r_\cA(z)$\textup{)}, and $5$
\textup{(}scalar diagonal $F_g^{\mathrm{diag}}$ and modular shadow
invariants\textup{)}; hypothesis package: Koszul chiral algebra $\cA$
with finite-type or weight-completed bar coalgebra; bar-intrinsic
MC element $\Theta_\cA$ \textup{(}Theorem~\ref{thm:mc2-bar-intrinsic}\textup{)};
modular reconstruction on the open factorization category for the
level-$5$ scalar/shadow projections.)
The modular invariants of a Koszul chiral algebra organize into
a hierarchy of increasing nonlinear depth:
\begin{equation}
\label{eq:full-hierarchy}
\underbrace{\kappa(\cA),\ \delta F_\bullet^{\mathrm{cross}}(\cA)}_{\text{scalar and mixed genus}}
\;\longrightarrow\;
\underbrace{\Delta_{\cA}}_{\text{spectral}}
\;\longrightarrow\;
\underbrace{[\mathfrak{C}_{\cA}],\;
\mathfrak{o}^{(4)}_{\cA},\;
[\mathfrak{Q}_{\cA}]}_{\text{nonlinear shadows}}
\;\longrightarrow\;
\underbrace{\Lambda_{P}(\mathfrak{Q}_{\cA})}_{\text{genus loop}}
\;\longrightarrow\;
\underbrace{\Theta_{\cA}}_{\text{bar-intrinsic limit}}.
\end{equation}
The projection order is functorial: the MC element determines the
finite obstruction tower, the finite tower determines its resonance
and spectral projections, and the diagonal trace gives the scalar
genus class. The shadow
obstruction tower
(Definition~\ref{def:shadow-postnikov-tower})
is the proved finite-order engine;
the all-degree limit $\Theta_{\cA}$
(Theorem~\ref{thm:mc2-bar-intrinsic}) is
the source of all of them.
In the tower language: the diagonal scalar projection
$\operatorname{tr}_{\mathrm{diag}}(\Theta_\cA)$ reads off
$\Theta_\cA^{\leq 2} = \kappa(\cA)$ (proved), while
$\operatorname{tr}_{\mathrm{cross}}(\Theta_\cA)$ records the mixed
boundary-graph projection of the same MC element (conditional through
Theorem~\ref{thm:multi-weight-genus-expansion}); the spectral
projection reads off $\Delta_\cA$ (proved at finite order); the
resonance tower reads off $\Theta_\cA^{\leq r}$ for
$r \geq 4$ (proved constructively; the all-degree limit exists
by the bar-intrinsic construction,
Theorem~\ref{thm:mc2-bar-intrinsic}).
\end{theorem}

\begin{proof}
The scalar characteristic $\kappa(\cA) = \operatorname{tr}(H_{\cA})$ is
the diagonal trace of the Hessian. The mixed scalar corrections are
the non-diagonal boundary contractions of the same graph expansion;
they vanish at genus~$1$ and on the scalar lane, and the first
standard nonzero witness is
$\delta F_2^{\mathrm{cross}}(\cW_3)=(c+204)/(16c)$.
The spectral characteristic $\Delta_{\cA}$
is the characteristic polynomial of the branching operator, which is
determined by the full Hessian. The cubic shadow
$[\mathfrak{C}_{\cA}]$ is a cocycle for $\nabla_{H_{\cA}}$
by the quartic master equation
(Theorem~\ref{thm:nms-shadow-master-equations}). The quartic shadow
$[\mathfrak{Q}_{\cA}]$ satisfies the driven equation
$\nabla_H\mathfrak{Q} + \frac12\{\mathfrak{C},\mathfrak{C}\} = 0$.
The genus loop $\Lambda_P(\mathfrak{Q}_{\cA})$ implements
non-separating clutching and produces the first tensorial modular
invariant beyond $\kappa(\cA)$
(Theorem~\ref{thm:nms-nonseparating-clutching-law}). Every level is
a projection of $\Theta_{\cA}$ under the Chern--Weil transform
\eqref{eq:modular-chern-weil-transform}.
\end{proof}

\begin{remark}[Beyond the $\hat{A}$-genus]
\label{rem:beyond-ahat-genus}
Two algebras with the same $\kappa(\cA)$ (hence identical $\hat{A}$-genera)
can be distinguished by their genus-$1$ Hessian correction
$\delta H^{(1)}_{\cA} = \Lambda_P(\mathfrak{Q}^{(0)}_{\cA})$.
For Virasoro,
$\delta H^{(1)}_{\mathrm{Vir}} = 120/[c^2(5c+22)] \cdot x^2$
and the loop ratio
$\rho^{(1)}_{\mathrm{Vir}} = 240/[c^3(5c+22)]$
is a nontrivial function of $c$ invisible to the scalar genus
expansion. Thus the modular homotopy type carries information
invisible to the $\hat{A}$-genus even when the diagonal scalar tower
agrees.
\end{remark}

% ==========================================
% ANALYTIC COMPLETION AND THE SEWING ENVELOPE
% ==========================================

\section{Analytic completion and the sewing envelope}
\label{sec:analytic-completion}
\index{sewing envelope|textbf}
\index{analytic completion|textbf}

The modular invariant hierarchy of the preceding section is algebraic:
it computes formal cohomological invariants of the bar complex, not
analytic data tied to Hilbert spaces, traces, or partition functions.
Convergent partition functions from local OPE data require an analytic
completion. Recent work of
Moriwaki~\cite{Moriwaki26a, Moriwaki26b} constructs conformally flat
factorization homology valued in $\operatorname{IndHilb}$ by left Kan
extension from conformally flat disk algebras, with sphere values
recovering CFT partition functions under positivity and continuity
assumptions; the Heisenberg case is made explicit via the Bergman
space~$\operatorname{Sym} A^2(D)$. Combined with the conformal
Osterwalder--Schrader axioms of
Adamo--Moriwaki--Tanimoto~\cite{AMT24} for unitary full VOAs, this
gives the analytic comparison surface for the algebraic genus tower.

\subsection{The sewing envelope}\label{subsec:sewing-envelope}
\index{sewing envelope!definition}

\begin{definition}[Sewing envelope]%
\label{def:sewing-envelope}%
Let $\cA_{\mathrm{alg}}$ be an algebraic chiral core with algebraic
amplitudes~$A_\Sigma^{\mathrm{alg}}$ for every finite conformally
flat bordism~$\Sigma$ with parametrized in/out collars. For every
matrix coefficient of every such amplitude, define a seminorm
on~$\cA_{\mathrm{alg}}$:
\[
p_{\Sigma,\xi,\eta}(a)
\;=\;
\bigl|\langle \eta,\,
A_\Sigma^{\mathrm{alg}}
(\xi_1 \otimes \cdots \otimes a \otimes \cdots \otimes \xi_m)
\rangle\bigr|.
\]
The \emph{sewing envelope}~$\cA^{\mathrm{sew}}$ is the Hausdorff
completion of~$\cA_{\mathrm{alg}}$ for the locally convex topology
generated by all seminorms~$p_{\Sigma,\xi,\eta}$.
\end{definition}

This is the minimal completion that remembers every local-to-global
operation at once.

\begin{proposition}[Universal property of the sewing envelope;
\ClaimStatusProvedHere]\label{prop:sewing-universal-property}%
\index{sewing envelope!universal property}%
If $E$ is any complete locally convex realization of\/
$\cA_{\mathrm{alg}}$ for which every algebraic amplitude extends
continuously, then the identity on~$\cA_{\mathrm{alg}}$ extends
uniquely to a continuous map
$\cA^{\mathrm{sew}} \to E$.
\end{proposition}

\begin{proof}
Every such realization makes each seminorm $p_{\Sigma,\xi,\eta}$
continuous. Hence the identity
$\cA_{\mathrm{alg}} \to E$ is continuous for the sewing topology.
By completeness of~$E$, it extends uniquely from the completion.
\end{proof}

\subsection{Hilbert--Schmidt sewing}\label{subsec:hs-sewing}
\index{Hilbert--Schmidt sewing|textbf}

Convergent partition functions require a quantitative condition on
the multiplication maps in addition to the formal gluing law.

\begin{definition}[HS-sewing condition]%
\label{def:hs-sewing}%
Let $H = \widehat{\bigoplus}_{n \geq 0} H_n$ be a graded pre-Hilbert
core with finite-dimensional energy sectors, and write the
pair-of-pants multiplication as sector maps
$m_{a,b}^c \colon H_a \,\widehat{\otimes}\, H_b \to H_c$.
Say $\cA$ satisfies \emph{HS-sewing at weight $q$} if
$0 < q < 1$ and the heat trace and pair-of-pants Hilbert--Schmidt
norms obey
\begin{equation}\label{eq:hs-heat-trace}
\sum_{n \geq 0} q^n \dim H_n < \infty,
\end{equation}
and
\begin{equation}\label{eq:hs-sewing}
\sum_{a,b,c \geq 0} q^{a+b+c}\,
\|m_{a,b}^c\|_{\mathrm{HS}}^2
\;<\; \infty.
\end{equation}
The collar transport and duality maps used to identify incoming and
outgoing boundary Hilbert spaces are required to be bounded on the
weighted completion below.  The algebra satisfies \emph{HS-sewing} if
it satisfies HS-sewing at some weight~$q$.
Define the \emph{weighted completion}
$H_q := \widehat{\bigoplus}_{n \geq 0} q^n H_n$.
\end{definition}

\begin{proposition}[Closed amplitudes are trace class;
\ClaimStatusProvedHere]\label{prop:hs-trace-class}%
\textbf{Type signature.} (Open quadrant; level~$5$ trace + clutching
analytic input on stable surfaces; hypothesis package: HS-sewing at
weight~$q$ \textup{(}Definition~\ref{def:hs-sewing}\textup{)}, bounded
collar transport on the weighted completion. Trace-class statements
on the closed amplitudes are level-$5$ analytic shadow consequences
of the open factorization category, not intrinsic closed-algebra
properties.)
Under HS-sewing at weight~$q$,
$m_q \colon H_q \,\widehat{\otimes}\, H_q \to H_q$ is
Hilbert--Schmidt. Every closed stable amplitude obtained by a finite
pair-of-pants sewing decomposition with collars controlled by~$q$ is
trace class. The genus-$1$ vacuum trace is trace class by the heat
trace condition~\eqref{eq:hs-heat-trace}.
\end{proposition}

\begin{proof}
$\|m_q\|_{\mathrm{HS}}^2 = \sum_{a,b,c} q^{a+b+c}
\|m_{a,b}^c\|_{\mathrm{HS}}^2 < \infty$ by~\eqref{eq:hs-sewing}.
For a closed stable surface of genus~$g\geq2$, a pair-of-pants
decomposition has $2g-2$ vertices and $3g-3$ internal sewing circles.
After choosing any internal circle, the amplitude factors through a
composition of two Hilbert--Schmidt maps and bounded collar
transports; such a composition is trace class. Separating degenerations
give tensor products of the same trace-class factors.  At genus~$1$
the closed vacuum amplitude is $\operatorname{Tr}(q^{L_0})$, finite by
\eqref{eq:hs-heat-trace}. The unstable genus-$0$ vacuum component is
not a trace-class assertion; it is the algebraic unit/counit sector.
\end{proof}

\begin{theorem}[General HS-sewing criterion; \ClaimStatusProvedHere]%
\label{thm:general-hs-sewing}%
\index{HS-sewing!general criterion|textbf}%
\textbf{Type signature.} (Open quadrant; level~$5$ analytic input for
trace + clutching on the open category; hypothesis package: positive
energy, subexponential sector growth, polynomial OPE growth, bounded
collar transport.)
Let\/ $\cA$ be a positive-energy chiral algebra with
\textup{(i)}~subexponential sector growth $\log\dim H_n = o(n)$
and \textup{(ii)}~polynomial OPE growth
$|C^{c,k}_{a,i;\,b,j}| \leq K(a+b+c+1)^N$, measured in orthonormal
sector bases, and bounded collar transport on the positive-energy
Hilbert realization. Then $\cA$ satisfies HS-sewing at every
$0 < q < 1$.
\end{theorem}

\begin{proof}
The heat trace condition follows from subexponential sector growth:
for fixed $0<q<1$ choose $\varepsilon>0$ with $qe^\varepsilon<1$ and
use $\dim H_n\leq C_\varepsilon e^{\varepsilon n}$ for all large~$n$.
Then $\sum_n q^n\dim H_n$ converges.

\[
\|m_{a,b}^c\|_{\mathrm{HS}}^2
\;\leq\;
\dim H_a \cdot \dim H_b \cdot \dim H_c
\cdot K^2(a{+}b{+}c{+}1)^{2N}.
\]
By the power-mean inequality $(x+y+z)^p \leq 3^{p-1}(x^p + y^p + z^p)$
applied with $p = 2N$ and $x = a{+}1$, $y = b{+}1$, $z = c{+}1$:
\[
(a{+}b{+}c{+}1)^{2N} \;\leq\; (a{+}1 + b{+}1 + c{+}1)^{2N}
\;\leq\; 3^{2N-1}\bigl((a{+}1)^{2N} + (b{+}1)^{2N} + (c{+}1)^{2N}\bigr).
\]
Substituting and using the symmetry of the three terms gives
$\sum_{a,b,c} q^{a+b+c}\|m_{a,b}^c\|_{\mathrm{HS}}^2
\leq 3^{2N}K^2
\bigl(\sum_n q^n\dim H_n\,(n{+}1)^{2N}\bigr)
\bigl(\sum_n q^n\dim H_n\bigr)^2$.
Fix $0<q<1$ and choose $\varepsilon>0$ with
$qe^\varepsilon<1$.  Hypothesis~(i) gives
$\dim H_n\leq C_\varepsilon e^{\varepsilon n}$ for all sufficiently
large~$n$.  Hence
$q^n\dim H_n(n+1)^{2N}\leq
C_\varepsilon(qe^\varepsilon)^n(n+1)^{2N}$, and the displayed
series converge.
\end{proof}

\begin{corollary}[Standard landscape; \ClaimStatusProvedElsewhere]%
\label{cor:hs-sewing-standard-landscape}%
\index{HS-sewing!standard families|textbf}%
\[
\begin{array}{c|ccc|c}
\cA & \dim H_n & |C| & N & \sum_n q^n\dim H_n
\\[2pt]\hline
\rule{0pt}{12pt}
\mathcal{H}_k & p(n) & \leq k & 1
 & q^{-1/24}/\eta(q) \\
V_k(\fg) & \leq Ce^{\alpha\sqrt{n}}
 & \leq K(n{+}1) & 1
 & \chi_{V_k}(q) \\
\mathrm{Vir}_c & p(n) & \leq K(n{+}1)^2 & 2
 & q^{(1-c)/24}/\eta(q) \\
\mathcal{W}_N & \leq Ce^{\alpha\sqrt{n}}
 & \leq K(n{+}1)^{2N} & 2N & \chi_{\mathcal{W}_N}(q) \\
V_\Lambda & |\Lambda^*/\Lambda|\cdot p(n)^r
 & \leq K \cdot r & r
 & \Theta_\Lambda(q)/\eta(q)^r
\end{array}
\]
\end{corollary}

\begin{remark}[Completed algebras]\label{rem:hs-sewing-completed}
\index{HS-sewing!completed algebras}
The families in Corollary~\ref{cor:hs-sewing-standard-landscape} are
all finitely strongly generated. For completed algebras such as
$\mathcal{W}_{1+\infty}$ and completed Yangians, the HS-sewing
condition requires additional verification: the polynomial OPE growth
bounds $|C^c_{a,b}| \leq K(a{+}b{+}c{+}1)^M$ used in
Theorem~\ref{thm:general-hs-sewing} must be checked in the completed
(inverse-limit) topology, where the OPE coefficients are formal
power series. This is the analytic sewing
input in the completed bar-cobar theorem
(Theorem~\ref{thm:completed-bar-cobar-strong}).
\end{remark}

\begin{corollary}[Diagonal positive sewing implies Gram positivity;
\ClaimStatusProvedHere]%
\label{cor:hs-implies-gram}%
\index{HS-sewing!Gram positivity}%
Assume HS-sewing and the diagonal weight-multiset expansion of
Definition~\ref{def:weight-multiset}. Then the connected sewing
amplitudes $a_\cA(N)$ are non-negative, define a
positive-semidefinite inner product on Dirichlet polynomials
\textup{(}Definition~\textup{\ref{def:sewing-inner-product}}\textup{)},
and $\mathscr{M}_\cA(\alpha) \succ 0$ for $\alpha \geq 2$
\textup{(}Theorem~\textup{\ref{thm:gram-positivity}}\textup{)}.  For
non-diagonal interacting sectors, positivity of the corrected
coefficients is an additional reflection-positivity input, not a
consequence of Hilbert--Schmidt convergence alone.
\end{corollary}

\begin{proof}
$a_\cA(N) = \sum_i\sum_{\substack{d \mid N \\ N/d \geq w_i}}
d^{-1} \geq 0$ for all $N \geq 1$: each summand is
a positive rational number. This is the input to
Theorem~\ref{thm:gram-positivity}.
\end{proof}

\begin{theorem}[Heisenberg: one-particle sewing;
\ClaimStatusProvedHere]%
\label{thm:heisenberg-one-particle-sewing}%
\index{Heisenberg!sewing|textbf}%
\index{Fredholm determinant!Heisenberg}%
\index{Bergman space!sewing operator}%
\textbf{Type signature.} (Open quadrant; level~$5$ analytic
specialisation of the trace + clutching reconstruction to class~G;
hypothesis package: rank-$r$ free Heisenberg core $\cH_{V,\kappa}$,
Bergman space
$A^2(D)$ on the disk, collar length $t = -\log q$ with
$0 < q < 1$.)
The Heisenberg pair-of-pants amplitude is the second
quantisation $\Gamma(R)$ of the one-particle Bergman
restriction $R \colon A^2(D_{\mathrm{out}}) \to
A^2(D_{\mathrm{in},1}) \otimes A^2(D_{\mathrm{in},2})$,
with $|R_{n,m}| \leq Cq^{(n+m)/2}$.
The sewing operator $T_q$ on the reduced Bergman space
$A^2_0(\partial D)$ (constant mode removed) has
eigenvalues~$q^n$ ($n \geq 1$) and is trace class:
$\|T_q\|_1 = q/(1-q)$.
The holomorphic oscillator part of the genus-$g$ chiral partition
function is the Fredholm determinant
\begin{equation}\label{eq:heisenberg-fredholm}
Z_{g,\mathrm{osc}}^{\mathrm{ch}}(\mathcal{H}_{V,\kappa};\,\Omega)
\;=\;
\det\nolimits_F(1-K_g)^{-r}.
\end{equation}
At genus~$1$:
$
Z_{1,\mathrm{osc}}^{\mathrm{ch}}=\eta(\tau)^{-r}.
$
The doubled non-chiral scalar determinant
$(\operatorname{Im}\tau)^{-r/2}|\eta(\tau)|^{-2r}$ is a different
object: it is obtained only after pairing the chiral determinant with
its anti-holomorphic conjugate and inserting the zero-mode measure.
\end{theorem}

\begin{proof}
Wick's theorem: all amplitudes factor into propagator pairings.
The one-particle restriction
$R_{n,m} = \frac{k}{2\pi i}\oint\oint z^{-n-1}w^{-m-1}
G_{\mathrm{pants}}(z,w)\,dz\,dw$
satisfies $|R_{n,m}| \leq Cq^{(n+m)/2}$ from the collar length
$t = -\log q$. Hence $R$ is Hilbert--Schmidt and $\Gamma(R)$
is bounded on the weighted Fock completion. For the actual sewing
trace one uses $T_q$, whose eigenvalues are $q^n$. Thus $T_q$ is trace
class, $\|T_q\|<1$, and bosonic second quantisation gives
\[
\operatorname{Tr}_{\operatorname{Sym}(A^2_0\otimes V)}\Gamma(T_q)
=\det\nolimits_F(1-T_q)^{-r}
=\prod_{n\ge1}(1-q^n)^{-r}.
\]
Multiplication by the vacuum-energy factor $q^{-r/24}$ gives
$\eta(\tau)^{-r}$. The Belavin--Knizhnik and Polyakov determinant
formula belongs to the doubled scalar measure; it is compatible with
the chiral determinant, but it is not identical to it.
\end{proof}

\begin{remark}[Fredholm--Epstein connection]\label{rem:heisenberg-bc-connection}
The Fredholm determinant $\det(1 - K_q) = \eta(q)/q^{1/24}$ connects
to the Benjamin--Chang constrained Epstein zeta
$\varepsilon^1_s(R{=}1) = 4\zeta(2s)$ via the primary-counting
function (Proposition~\ref{prop:benjamin-chang-bridge}).
\end{remark}

\subsection{Analytic bar coalgebra and Koszul pairs}%
\label{subsec:analytic-bar-koszul}%
\index{bar coalgebra!analytic}

\begin{definition}[Analytic bar coalgebra]%
\label{def:analytic-bar-coalgebra}%
For an $\mathsf{E}_1$-type chiral core~$\cA_{\mathrm{alg}}$, define
the $n$-th completed bar piece by closing the algebraic bar sector
inside the sewing envelope:
\[
\barB_n^{\mathrm{an}}(\cA)
\;=\;
\overline{
\cA_{\mathrm{alg}}^{\otimes n}
\otimes \Omega^*_{\log}(C_n(X))
}^{\,\|\cdot\|_{\mathrm{graph}(d_{\mathrm{bar}})}},
\]
where the closure is taken in the graph norm of the closable
residue operator~$d_{\mathrm{bar}} = \sum_{i<j}
\operatorname{Res}_{z_i = z_j} \circ Y_{ij}$.
\end{definition}

This is the analytic closure of the algebraic bar construction
(Definition~\ref{def:bar-differential-complete}), not a different
construction.

\begin{proposition}[Finite-window conilpotency and completed pro-conilpotency;
\ClaimStatusProvedHere]\label{prop:analytic-conilpotency}%
Let
$C=\varprojlim_N C_{\leq N}$ be a coaugmented completed coalgebra
with $C_{\leq N}:=\bigoplus_{0\leq n\leq N}C_n$ and
$C_0=\mathbb C$. Suppose the reduced coproduct preserves total degree:
$\bar{\Delta}(C_n) \subset
\widehat{\bigoplus}_{i+j=n,\; i,j>0} C_i \,\widehat{\otimes}\, C_j$.
Then each finite window $C_{\leq N}$ is conilpotent, and the
completed coalgebra $C$ is pro-conilpotent:
for every $x\in C$ and every finite window $C_{\leq N}$, the image of
$\bar{\Delta}^{(M)}(x)$ in $C_{\leq N}^{\widehat{\otimes}M}$
vanishes for all $M>N$.
\end{proposition}

\begin{proof}
$\bar{\Delta}(C_n) \subset
\bigoplus_{i+j=n,\, i,j > 0} C_i \,\widehat{\otimes}\, C_j$.
Iterating: $\bar{\Delta}^{(N)}(C_n)$ lands in
$\bigoplus_{i_1+\cdots+i_N=n,\, i_k > 0}
C_{i_1} \,\widehat{\otimes}\, \cdots \,\widehat{\otimes}\, C_{i_N}$.
For $N > n$, the sum $i_1 + \cdots + i_N = n$ with each
$i_k \geq 1$ has no solutions.
Hence $\bar{\Delta}^{(N)}(C_n) = 0$ for $N > n$.

Thus $\bar{\Delta}^{(M)}$ vanishes on $C_{\leq N}$ for $M>N$.
For $x\in C=\varprojlim_N C_{\leq N}$, write $x_{\leq N}$ for its
image in the finite window. The compatibility of $\bar{\Delta}$ with
the inverse system gives
\[
\bigl(\bar{\Delta}^{(M)}x\bigr)_{\leq N}
=\bar{\Delta}^{(M)}(x_{\leq N})=0
\qquad (M>N).
\]
This is exactly convergence of iterated reduced coproducts to zero in
the completion topology. Literal elementwise conilpotency is asserted
only inside finite windows; the completed bar object is
pro-conilpotent.

For the analytic bar coalgebra, the same proof applies when the graph
norm respects the degree filtration and $d_{\mathrm{bar}}$ preserves
that filtration (Definition~\ref{def:analytic-bar-coalgebra}); bounded
OPE pole order gives this condition for the standard-landscape
families.
\end{proof}

\begin{definition}[Analytic Koszul pair]%
\label{def:analytic-koszul-pair}%
\index{Koszul pair!analytic}%
An \emph{analytic Koszul pair} $(\cA, \cA^!)$ consists of:
\begin{enumerate}[label=\textup{(\roman*)}]
\item a sewing envelope $\cA^{\mathrm{sew}}$
 (Definition~\ref{def:sewing-envelope}),
\item an analytic bar coalgebra $\barB^{\mathrm{an}}(\cA)$
 (Definition~\ref{def:analytic-bar-coalgebra}),
\item a genus-zero equivalence
 $D^{\mathrm{an}}_{\mathrm{compl}}(\cA\text{-mod})
 \simeq
 D^{\mathrm{an}}_{\mathrm{conil}}
 (\barB^{\mathrm{an}}(\cA)\text{-comod})$,
\item under finite-type dualizability, an identification of
 conilpotent analytic comodules with complete
 $\cA^!$-modules.
\end{enumerate}
\end{definition}

\begin{remark}[Curvature forces coderived passage]%
\label{rem:curvature-coderived-analytic}%
\index{coderived category!analytic motivation}%
At genus~$g$, if the completed bar object carries a curvature term
$\mcurv{g} \neq 0$, then ordinary cohomology fails because the
differential satisfies $d^2 = \mcurv{g} \cdot (-)$, and honest
cohomology stops being functorial under sewing. The correct
receptacle is the coderived/contraderived analytic category
(\S\ref{subsec:coderived-ran}), not an ordinary derived category.
This is the mechanism by which higher-genus local-to-global data
remain nontrivial after completion.
\end{remark}

\subsection{Analytic shadows and the shadow partition function}%
\label{subsec:analytic-shadows}%
\index{shadow!analytic}

\begin{definition}[Analytic shadow]%
\label{def:analytic-shadow}%
For each genus~$g$, define the \emph{analytic shadow}
\[
Q_g^{\mathrm{an}}(\cA)
\;:=\;
D^{\mathrm{co}}\!\bigl(
\barB_g^{\mathrm{an}}(\cA)\text{-comod}\bigr).
\]
When $\mcurv{g} = 0$, this reduces to the cohomological shadow
$\operatorname{Sh}_g(\cA) = H^*\Gamma(\barB_g^{\mathrm{an}}(\cA))$.
A \emph{complementarity operator} is an involution
$\sigma_g \colon Q_g^{\mathrm{an}}(\cA) \to Q_g^{\mathrm{an}}(\cA)$
induced by duality plus sewing reversal.
\end{definition}

\begin{definition}[Shadow partition function]%
\label{def:shadow-partition-function}%
\index{partition function!shadow}%
When $q^{L_0}\sigma_g$ is trace class on $Q_g^{\mathrm{an}}(\cA)$,
define the \emph{shadow partition function}
\begin{equation}\label{eq:shadow-partition-function}
Z_g^{\mathrm{sh}}(\cA; \tau)
\;=\;
\operatorname{Str}_{Q_g^{\mathrm{an}}(\cA)}
\!\bigl(\sigma_g\, q^{L_0 - c/24}\bigr).
\end{equation}
By Proposition~\ref{prop:hs-trace-class}, HS-sewing is the hypothesis
that makes this analytic trace converge for $|q|$ sufficiently below
$1$; modular or meromorphic continuation is a separate input.
\end{definition}

\begin{conjecture}[Centers survive analytic Koszul duality;
\ClaimStatusConjectured]\label{conj:analytic-center-invariance}%
Assume an analytic Koszul equivalence
$\Phi \colon D^{\mathrm{an}}_{\mathrm{compl}}(\cA)
\xrightarrow{\,\sim\,}
D^{\mathrm{an}}_{\mathrm{compl}}(\cA^!)$
sending vacuum to vacuum. Then
\[
Z^{\mathrm{an}}(\cA) \;\cong\; Z^{\mathrm{an}}(\cA^!),
\qquad
Z^{\mathrm{an}}(\cA) := \operatorname{End}(\mathbf{1}_\cA).
\]
\end{conjecture}

\begin{remark}
The heuristic support: equivalences preserve endomorphism algebras of
distinguished objects. A full proof requires the analytic Koszul
equivalence to be established, which is the content of
Conjecture~\ref{conj:analytic-realization}.
\end{remark}

Shadow and complementarity invariants are therefore completed
objects: the bare algebraic presentation supplies generators and OPE
relations, while the sewing envelope supplies the topology in which
traces and coderived shadows exist.

\subsection{The completed analytic chain}%
\label{subsec:platonic-chain}

\begin{remark}[Completed analytic form]%
\label{rem:platonic-ideal-form}%
\index{analytic completion!completed form}%
The analytic form of a chiral algebra is the chain
\begin{equation}\label{eq:platonic-chain}
\underbrace{\cA_{\mathrm{alg}}}_{\text{algebraic core}}
\;\subset\;
\underbrace{\cA^{\mathrm{sew}}}_{\text{sewing envelope}}
\;\rightsquigarrow\;
\underbrace{F_\cA \in
\operatorname{IndHilb}}_{\text{Hilbert factorization theory}}
\;\rightsquigarrow\;
\underbrace{Q_\bullet^{\mathrm{an}}(\cA)}_{\text{coderived shadow
invariants}},
\end{equation}
where $F_\cA$ denotes the $\operatorname{IndHilb}$-valued conformally
flat factorization theory obtained by left Kan extension
\`a~la~Moriwaki~\cite{Moriwaki26a}. The algebraic core determines the
local OPE; the sewing envelope, the boundary/defect module theory
$M_\partial^{\mathrm{an}}(\cA)$, and the coderived genus tower
$Q_\bullet^{\mathrm{an}}(\cA)$ determine the global analytic shadow.

If the formal shadow invariants of
\S\ref{subsec:nonlinear-modular-hierarchy} have a
nonlinear modular shadow, it appears after completion, on
completed factorization homology, boundary Hilbert spaces, and
analytic sewing data, where formal shadow invariants can interact
with unitarity and partition-function data.
\end{remark}

\begin{conjecture}[Analytic realization criterion;
\ClaimStatusConjectured]\label{conj:analytic-realization}%
\index{analytic realization criterion}%
Let $V$ be a unitary full VOA satisfying:
\begin{enumerate}[label=\textup{(\roman*)}]
\item its quasi-primary correlators satisfy conformal OS axioms
 \textup{(}in the sense of Adamo--Moriwaki--Tanimoto
 \cite{AMT24}\textup{)},
\item its graded pieces have polynomial spectral growth,
\item its intertwining OPE maps satisfy HS-sewing
 \textup{(}Definition~\textup{\ref{def:hs-sewing}}\textup{)}
 after exponential weight damping.
\end{enumerate}
Then $V$ admits a sewing envelope
$V_{\mathrm{alg}} \to V^{\mathrm{sew}} \in
\operatorname{IndHilb}$, a conformally flat 2-disk algebra~$F_V$, a
genus-zero analytic Koszul duality package, and a higher-genus
coderived shadow $Q_\bullet^{\mathrm{an}}(V)$.
\end{conjecture}

\begin{remark}[Heisenberg core and analytic extension]%
\label{rem:analytic-execution-order}%
\index{analytic completion!Heisenberg core}%
The proved analytic inputs above are the sewing-envelope universal
property, the HS trace-class criterion, finite-window
pro-conilpotency, and the Heisenberg one-particle computation. The
remaining conjectural step is the realization criterion
(Conjecture~\ref{conj:analytic-realization}) for unitary full VOAs
covered by the conformal OS theorem. Lattice theories add charge
sectors through theta data; higher-genus coderived shadows then carry
the complementarity involution.
\end{remark}

\begin{definition}[Analytic boundary condition]%
\label{def:analytic-boundary-condition}%
\index{boundary condition!analytic}%
An \emph{analytic boundary condition} for~$\cA$ is a complete
module~$M$ equipped with continuous half-disk actions
$\cA(U) \,\widehat{\otimes}\, M(I) \to M(I')$
compatible with collar sewing and reflection.
\end{definition}

\begin{conjecture}[Boundary bar duality;
\ClaimStatusConjectured]\label{conj:boundary-bar-duality}%
\index{boundary condition!bar duality}%
For an analytically Koszul~$\cA$
\textup{(}Definition~\textup{\ref{def:analytic-koszul-pair}}\textup{)}:
\begin{enumerate}[label=\textup{(\roman*)}]
\item at genus~$0$, completed boundary modules
 \textup{(}Definition~\textup{\ref{def:analytic-boundary-condition}}\textup{)}
 are equivalent to analytic comodules over
 $\barB^{\mathrm{an}}(\cA)$;
\item at curved genus, the correct objects are analytic contramodules
 over the curved dual.
\end{enumerate}
This is where branes, open strings, and boundary chiral algebras
should live in the analytic framework.
\end{conjecture}

\subsection{The Dirichlet--sewing lift}%
\label{subsec:arithmetic-sewing-face}%
\label{subsec:dirichlet-sewing-lift}%
\index{Dirichlet--sewing lift|textbf}%
\index{Euler--Koszul|textbf}

\begin{definition}[Connected genus-$1$ free energy]%
\label{def:connected-genus1-free-energy}%
\index{free energy!connected genus-1}%
Let $\cA$ satisfy HS-sewing
(Definition~\ref{def:hs-sewing}) with genus-$1$ sewing
operator~$K_q(\cA)$.
The \emph{connected genus-$1$ free energy}
\begin{equation}\label{eq:connected-free-energy}
F_\cA^{\mathrm{conn}}(q)
\;:=\;
-\log\det(1 - K_q(\cA))
\;=\;
\sum_{N=1}^\infty a_\cA(N)\, q^N
\end{equation}
is the genus-$1$, degree-$0$ component of
$\log Z_\cA$: the connected vacuum amplitude on the torus.
For the Heisenberg algebra
(Theorem~\ref{thm:heisenberg-one-particle-sewing}),
$F_\cH^{\mathrm{conn}} = -\log\prod_{n \geq 1}(1-q^n) =
\sum_{n \geq 1}\sum_{r \geq 1} q^{nr}/r$.
\end{definition}

\begin{definition}[Dirichlet--sewing lift]%
\label{def:dirichlet-sewing-lift}%
\index{Dirichlet--sewing lift|textbf}%
\index{Dirichlet series!sewing}%
The \emph{Dirichlet--sewing lift} of~$\cA$ is
\begin{equation}\label{eq:dirichlet-sewing-lift}
S_\cA(u)
\;:=\;
\sum_{N=1}^\infty a_\cA(N)\, N^{-u}
\;=\;
\frac{(2\pi)^u}{\Gamma(u)}
\int_0^\infty
F_\cA^{\mathrm{conn}}(e^{-2\pi y})\,y^{u-1}\,dy,
\end{equation}
The equality is first asserted in the common domain of absolute
convergence. In the diagonal weight-multiset sector of
Definition~\ref{def:weight-multiset} this domain contains
$\operatorname{Re}(u)>1$; HS-sewing supplies the trace-class
$q$-amplitude, but does not by itself determine the exact Dirichlet
abscissa.
\end{definition}

\begin{definition}[Weight multiset]%
\label{def:weight-multiset}%
\index{weight multiset|textbf}%
Assume the genus-$1$ sewing operator is diagonal on a one-particle
mode basis with one oscillator string for each strong generator. The
\emph{weight multiset}
$W(\cA) = \{w_1, \ldots, w_r\}$
is the multiset of conformal weights of the strong generators
of~$\cA$. Its diagonal connected free energy is the
weight-multiset expansion
\begin{equation}\label{eq:weight-multiset-expansion}
F_\cA^{\mathrm{conn}}(q)
\;=\;
\sum_{i=1}^r\,
\sum_{m \geq w_i}\,
\sum_{k \geq 1}
\frac{q^{mk}}{k}\,,
\end{equation}
with Dirichlet coefficients
$a_\cA(N) = \sum_{i=1}^r
\sum_{\substack{d \mid N,\; N/d \geq w_i}}d^{-1}$.
\end{definition}

\begin{theorem}[\ClaimStatusProvedHere]%
\label{thm:dirichlet-weight-formula}%
\index{Dirichlet--sewing lift!weight formula}%
Then
\begin{equation}\label{eq:dirichlet-weight-formula}
S_\cA(u)
\;=\;
\zeta(u{+}1)
\sum_{i=1}^r
\bigl(\zeta(u) - H_{w_i - 1}(u)\bigr),
\qquad
H_n(u) := \sum_{j=1}^n j^{-u}.
\end{equation}
In particular:
\begin{equation}\label{eq:dirichlet-standard}
\renewcommand{\arraystretch}{1.3}
\begin{array}{r@{\;:\;}l@{\qquad}l}
\cH & S_\cH(u) = \zeta(u)\,\zeta(u{+}1)
 & = \prod_p (1{-}p^{-u})^{-1}(1{-}p^{-u-1})^{-1},
\\
\beta\gamma & S_{\beta\gamma}(u) = 2\,\zeta(u)\,\zeta(u{+}1),
\\
V_k(\fg) & S_{V_k(\fg)}(u) = \dim(\fg)\cdot\zeta(u)\,\zeta(u{+}1),
\\
\mathrm{Vir}_c & S_{\mathrm{Vir}}(u) = \zeta(u{+}1)\bigl(\zeta(u) - 1\bigr),
\\
\cW_N & S_{\cW_N}(u)
 = \zeta(u{+}1)\Bigl((N{-}1)\zeta(u) -
 \sum_{j=1}^{N-1} H_j(u)\Bigr).
\end{array}
\end{equation}
\end{theorem}

\begin{proof}
A strong generator of conformal weight~$w$ contributes
modes at levels $m \geq w$ to the sewing operator,
each with multiplicity~$1$. Hence
\[
F_\cA^{\mathrm{conn}}(q)
= \sum_{i=1}^r \sum_{m \geq w_i}\sum_{k \geq 1}
 \frac{q^{mk}}{k},
\qquad
a_\cA(N) = \sum_{i=1}^r
 \sum_{\substack{d \mid N \\ N/d \,\geq\, w_i}}
 \frac{1}{d}.
\]
Taking the Dirichlet series and exchanging:
$S_\cA(u) = \sum_i \sum_{m \geq w_i} m^{-u}\,\zeta(u{+}1)
= \zeta(u{+}1)\sum_i(\zeta(u) - H_{w_i-1}(u))$.
For Heisenberg: $H_0 = 0$,
$a_\cH(N) = \sum_{d \mid N} d^{-1} = \sigma_{-1}(N)$
(multiplicative), giving the Euler product.
For $V_k(\fg)$: $W = \{1^{\dim\fg}\}$
($J^a$ for $a = 1,\ldots,\dim\fg$;
Sugawara is composite), each summand $= \zeta(u)$.
\end{proof}

\begin{corollary}[\ClaimStatusProvedHere]%
\label{cor:virasoro-mode-removal}%
\index{Virasoro!mode removal}%
\index{Drinfeld--Sokolov!sewing subtraction}%
The Virasoro lift factors as mode removal from the
Heisenberg Euler product:
\begin{equation}\label{eq:virasoro-mode-removal}
S_{\mathrm{Vir}}(u)
\;=\;
\zeta(u)\,\zeta(u{+}1)
\bigl(1 - \zeta(u)^{-1}\bigr)
\;=\;
S_\cH(u)\cdot D_{\mathrm{Vir}}(u).
\end{equation}
More generally, for $\cW_N$:
\begin{equation}\label{eq:wn-alternative}
S_{\cW_N}(u)
\;=\;
\zeta(u{+}1)
\Bigl(
(N{-}1)\,\zeta(u)
- N\,H_{N-1}(u)
+ H_{N-1}(u{-}1)
\Bigr),
\end{equation}
where $H_n(u{-}1) = \sum_{j=1}^n j^{-(u-1)}$.
\end{corollary}

\begin{proof}
\eqref{eq:virasoro-mode-removal}:
$\zeta(u) - 1 = \zeta(u)(1 - \zeta(u)^{-1})$.
\eqref{eq:wn-alternative}: from
$\sum_{j=1}^{N-1}H_j(u)
= \sum_{j=1}^{N-1}\sum_{k=1}^{j}k^{-u}
= \sum_{k=1}^{N-1}(N{-}k)\,k^{-u}
= N\,H_{N-1}(u) - H_{N-1}(u{-}1)$.
\end{proof}

\begin{definition}[Euler--Koszul defect and class]%
\label{def:euler-koszul-defect}%
\label{def:euler-koszul-class}%
\index{Euler--Koszul defect|textbf}%
\index{Euler--Koszul!class|textbf}%
The \emph{Euler--Koszul defect}
\begin{equation}\label{eq:euler-koszul-defect}
D_\cA(u)
\;:=\;
\frac{S_\cA(u)}{|W(\cA)|\cdot\zeta(u)\,\zeta(u{+}1)}
\;=\;
\frac{1}{r}\sum_{i=1}^r
\Bigl(1 - \frac{H_{w_i-1}(u)}{\zeta(u)}\Bigr)
\end{equation}
is a meromorphic function of~$u$ with poles only at zeros
of~$\zeta$.
The \emph{Euler--Koszul class}
$\operatorname{ek}(\cA) := \max_i(w_i - 1)$
is the maximal generator weight excess.
\end{definition}

\begin{theorem}[\ClaimStatusProvedHere]%
\label{thm:euler-koszul-criterion}%
\label{thm:euler-koszul-standard}%
\index{Euler--Koszul!criterion}%
\begin{enumerate}[label=\textup{(\roman*)}]
\item
$D_\cA \equiv 1$ if and only if $w_i = 1$ for all~$i$.
\item
$D_{\mathrm{Vir}}(u) = 1 - 1/\zeta(u)$; \quad
$D_{\cW_N}(u)
= 1 - \frac{1}{(N{-}1)\zeta(u)}
\sum_{j=1}^{N-1} H_j(u)$.
\item
Principal DS reduction $V_k(\fg) \leadsto \cW_k(\fg)$
replaces $\dim\fg$ weight-$1$ generators by $r$ generators
of weights $d_i + 1$ \textup{(}$d_i$ the exponents
of~$\fg$\textup{)}:
\begin{equation}\label{eq:ds-sewing-subtraction}
S_{\cW_k(\fg)}(u)
\;=\;
\zeta(u{+}1)\sum_{i=1}^{r}
\bigl(\zeta(u) - H_{d_i}(u)\bigr).
\end{equation}
\end{enumerate}
\end{theorem}

\begin{proof}
(i)~$w_i = 1 \Leftrightarrow H_{w_i-1} = 0
\Leftrightarrow$ each summand in~\eqref{eq:euler-koszul-defect}
equals~$1$.
(ii)~Substitute $W = \{2\}$ and $\{2,\ldots,N\}$.
(iii)~$\cW_k(\fg)$ is strongly generated by $W$-currents
$W^{(d_i+1)}$ of conformal weight~$d_i + 1$;
substitute into Theorem~\ref{thm:dirichlet-weight-formula}.
\end{proof}

\begin{computation}[Euler--Koszul defect table for the standard landscape;
\ClaimStatusProvedHere]%
\label{comp:euler-koszul-defect-table}%
\index{Euler--Koszul defect!standard families|textbf}%
Substitute~\eqref{eq:dirichlet-standard}
into~\eqref{eq:euler-koszul-defect}:
\begin{center}
\renewcommand{\arraystretch}{1.35}
\begin{tabular}{lccccc}
\toprule
\emph{Family} & $W(\cA)$ & $S_\cA(u)$ & $D_\cA(u)$ & \emph{Tier}
 & $r_D$ \\
\midrule
$\cH_c$ & $\{1\}$ & $\zeta\,\zeta'$
 & $1$ & $1$ & $0$ \\
$\beta\gamma$ & $\{1,1\}$ & $2\,\zeta\,\zeta'$
 & $1$ & $1$ & $0$ \\
$V_k(\mathfrak{sl}_2)$ & $\{1^3\}$ & $3\,\zeta\,\zeta'$
 & $1$ & $1$ & $0$ \\
$V_k(\mathfrak{sl}_N)$ & $\{1^{N^2-1}\}$
 & $(N^2{-}1)\,\zeta\,\zeta'$
 & $1$ & $1$ & $0$ \\[3pt]
$\mathrm{Vir}_c$ & $\{2\}$
 & $\zeta'(\zeta{-}1)$
 & $1 - \dfrac{1}{\zeta(u)}$ & $2$ & $1$ \\[6pt]
$\cW_3$ & $\{2,3\}$
 & $\zeta'\bigl(2\zeta-1-H_2\bigr)$
 & $1 - \dfrac{2+2^{-u}}{2\,\zeta(u)}$ & $2$ & $2$ \\[6pt]
$\cW_N$ & $\{2,\ldots,N\}$
 & $\zeta'\bigl((N{-}1)\zeta-\sum H_j\bigr)$
 & $1 - \dfrac{1}{(N{-}1)\,\zeta(u)}
 \displaystyle\sum_{j=1}^{N-1}H_j(u)$
 & $2$ & $N{-}1$ \\
\bottomrule
\end{tabular}
\end{center}
Here $\zeta = \zeta(u)$, $\zeta' = \zeta(u{+}1)$,
$H_n = H_n(u)$, and
\emph{$r_D$} is the defect rank: the number of finite harmonic
subtraction summands in $1-D_\cA(u)$ after collecting the
weight-multiset contribution. It is not a polynomial degree in
$1/\zeta(u)$.
The defining formulas are:
\begin{align}
D_{\cH}(u) &= 1,
\label{eq:defect-heisenberg}
\\
D_{\mathrm{Vir}}(u) &= 1 - \frac{1}{\zeta(u)},
\label{eq:defect-virasoro}
\\
D_{\cW_3}(u) &= 1 - \frac{1 + H_2(u)}{2\,\zeta(u)}
= 1 - \frac{2 + 2^{-u}}{2\,\zeta(u)},
\label{eq:defect-w3}
\\
D_{\cW_N}(u) &= 1 -
\frac{1}{(N{-}1)\,\zeta(u)}\sum_{j=1}^{N-1}H_j(u).
\label{eq:defect-wn}
\end{align}
For the $\cW_3$ defect, $H_1(u) = 1$ and
$H_2(u) = 1 + 2^{-u}$, giving
$D_{\cW_3}(u) = 1 - (2 + 2^{-u})/(2\zeta(u))$,
which has defect rank~$2$: the weight-$2$ generator removes mode~$1$,
and the weight-$3$ generator removes modes~$1$ and~$2$.

\smallskip
\noindent\textbf{DS tier transition.}
Principal DS reduction
$V_k(\mathfrak{sl}_N) \leadsto \cW_k(\mathfrak{sl}_N)$
sends Tier~$1$ to Tier~$2$:
the $N^2 - 1$ weight-$1$ generators~$\{1^{N^2-1}\}$
are replaced by $N - 1$ generators of weights
$\{2, 3, \ldots, N\}$
(Casimir degrees of~$\mathfrak{sl}_N$).
The arithmetic defect polynomial
$P_{\cW_N}(u) := 1 - D_{\cW_N}(u)
= \frac{1}{(N{-}1)\,\zeta(u)}\sum_{j=1}^{N-1} H_j(u)$
has defect rank~$N - 1 = \operatorname{rk}(\mathfrak{sl}_N)$,
one harmonic subtraction for each Casimir degree.
In general, for any simple~$\fg$ of rank~$r$
with exponents $d_1,\ldots,d_r$:
\begin{equation}\label{eq:ds-defect-degree}
r_D(\cW_k(\fg)) = r = \operatorname{rk}(\fg),
\qquad
P_{\cW_k(\fg)}(u) =
\frac{1}{r\,\zeta(u)}\sum_{i=1}^{r} H_{d_i}(u).
\end{equation}
Each generator of conformal weight $d_i + 1$ contributes
$H_{d_i}(u)/\zeta(u)$
to the defect: it removes modes $m = 1, \ldots, d_i$
from the Euler product. The maximal removed mode is
$\operatorname{ek}(\cW_k(\fg))=\max_i d_i$.
\end{computation}

\begin{proposition}[\ClaimStatusProvedHere]%
\label{prop:zeta-zeros-defect-derivative}%
\index{Riemann zeta function!zeros in defect derivative}%
\begin{equation}\label{eq:defect-derivative}
\frac{D'_{\mathrm{Vir}}(u)}{D_{\mathrm{Vir}}(u)}
\;=\;
\frac{\zeta'(u)}{\zeta(u)\bigl(\zeta(u) - 1\bigr)}\,,
\end{equation}
with poles at every nontrivial zero of~$\zeta$.
\end{proposition}

\begin{proof}
$D = 1 - \zeta^{-1}$;
$D' = \zeta'/\zeta^2$;
$D'/D = \zeta'/(\zeta(\zeta-1))$.
At a zero $\rho$ of multiplicity~$m$, this quotient has a simple pole
with residue $-m$.
\end{proof}

\subsection{Prime-side Li coefficients}%
\label{subsec:prime-side-li}%
\index{Li coefficients!prime-side|textbf}

\begin{definition}[Prime-side Li coefficients]%
\label{def:prime-side-li}%
\index{Dirichlet--sewing lift!Li coefficients}%
Set $\Xi_\cA(u) := (u{-}1)\,S_\cA(u)$. In the weight-multiset cases
below this removes the pole at $u=1$, so $\Xi_\cA$ is holomorphic and
nonzero in a neighbourhood of~$u=1$. It need not be entire: the
factor $\zeta(u{+}1)$ has its pole at $u=0$.
The \emph{prime-side Li coefficients} are
\begin{equation}\label{eq:prime-side-li}
\tilde\lambda_n(\cA)
\;:=\;
\frac{1}{(n{-}1)!}
\frac{d^n}{du^n}
\Bigl[u^{n-1}\log\Xi_\cA(u)\Bigr]_{u=1}\,.
\end{equation}
\end{definition}

\begin{theorem}[\ClaimStatusProvedHere]%
\label{thm:li-closed-form}%
\index{Li coefficients!closed form}%
\begin{equation}\label{eq:li1-heisenberg}
\tilde\lambda_1(\cH) = \gamma + \frac{\zeta'(2)}{\zeta(2)}\,,
\qquad
\tilde\lambda_1(\cW_N)
= \frac{\zeta'(2)}{\zeta(2)} + \gamma + 1
 - \frac{N}{N{-}1}\,H_{N-1}\,.
\end{equation}
\end{theorem}

\begin{proof}
$\Xi_\cH(u) = (u{-}1)\zeta(u)\cdot\zeta(u{+}1)$.
The Laurent expansion $(u{-}1)\zeta(u) =
1 + \gamma(u{-}1) + \gamma_1(u{-}1)^2 + \cdots$
(with $\gamma_k$ the $k$-th Stieltjes constant)
gives $\Xi_\cH(1) = \zeta(2)$ and
$(\log\Xi_\cH)'(1)
= \gamma + \zeta'(2)/\zeta(2)$.
For~$\cW_N$:
$\Xi_{\cW_N}(1) = (N{-}1)\zeta(2)$
(each term $(u{-}1)H_j(u) \to 0$).
Differentiating:
$(d/du)[(u{-}1)H_j(u)]_{u=1} = H_j$,
so
$(\log\Xi_{\cW_N})'(1)
= \gamma + \zeta'(2)/\zeta(2)
+ 1/(N{-}1) \cdot
\bigl((N{-}1) - \sum_{j=1}^{N-1}H_j\bigr)$.
Rearranging gives the formula.
\end{proof}

\begin{theorem}[First Li coefficient and finite sign computation;
\ClaimStatusProvedHere]%
\label{thm:li-asymptotics}%
\index{Li coefficients!asymptotics}%
\begin{enumerate}[label=\textup{(\roman*)}]
\item
$\displaystyle
\lim_{N \to \infty}
\bigl(\tilde\lambda_1(\cW_N) + \log N\bigr)
= \frac{\zeta'(2)}{\zeta(2)} + 1
\approx 0.4300$.
\item
$\tilde\lambda_1(\cW_N)<0$ for every $N\geq2$.
\item
Direct 50-digit evaluation of the defining derivatives gives
$\tilde\lambda_n(\cH) > 0$ for $1 \leq n \leq 6$,
$\tilde\lambda_7(\cH) < 0$, and
$\tilde\lambda_n(\mathrm{Vir})<0$,
$\tilde\lambda_n(\cW_3)<0$,
$\tilde\lambda_n(\cW_4)<0$ for $1\leq n\leq20$.
\end{enumerate}
\end{theorem}

\begin{proof}
(i)~Write $\frac{N}{N{-}1}H_{N-1}
= (1 + \tfrac{1}{N{-}1})
\bigl(\log(N{-}1) + \gamma + \tfrac{1}{2(N{-}1)} + O(N^{-2})\bigr)
= \log(N{-}1) + \gamma + O(1/N)
= \log N + \gamma + O(1/N)$,
using $\log(N{-}1) = \log N - 1/N + O(N^{-2})$.
Substituting into~\eqref{eq:li1-heisenberg}:
$\tilde\lambda_1(\cW_N) + \log N
= \zeta'(2)/\zeta(2) + \gamma + 1 - \gamma - O(1/N)
= \zeta'(2)/\zeta(2) + 1 + O(1/N)$.
(ii)~Since
$\zeta'(2)/\zeta(2)=-\sum_{m\geq2}(\log m)m^{-2}/\zeta(2)<0$
and $\gamma<1$, the case $N=2$ gives
$\tilde\lambda_1(\cW_2)=\zeta'(2)/\zeta(2)+\gamma-1<0$.
For $N=3,4$ the factor is respectively
$\frac32(1+\frac12)>2$ and
$\frac43(1+\frac12+\frac13)>2$; for $N\geq5$ one has
$H_{N-1}\geq H_4=25/12>2$, hence
$\frac{N}{N-1}H_{N-1}>2$.  Thus the same upper bound applies.
(iii)~The finite sign range is a numerical evaluation of
Definition~\ref{def:prime-side-li}. No all-order sign theorem for
$n>20$ is asserted here.
\end{proof}

\begin{remark}[Numerical values and the two Li theories]%
\label{rem:li-numerical-two-theories}%
\index{Li coefficients!numerical values}%
\index{Li coefficients!two theories}%
The first two prime-side Li coefficients for the standard
families (50-digit computation):
\begin{center}
\small
\renewcommand{\arraystretch}{1.2}
\begin{tabular}{lrr}
\toprule
\emph{Family} & $\tilde\lambda_1$ & $\tilde\lambda_2$ \\
\midrule
$\cH$ & $+0.0072546718$ & $+0.7114449447$ \\
$\mathrm{Vir} = \cW_2$ & $-0.9927453282$ & $-1.1341237255$ \\
$\cW_3$ & $-1.2427453282$ & $-1.5614423027$ \\
$\cW_4$ & $-1.4371897726$ & $-1.8901175919$ \\
\bottomrule
\end{tabular}
\end{center}
The displayed finite sign pattern is sharp in the computed range:
Heisenberg is weakly positive at low order and has its first computed
sign change at $n = 7$.  Virasoro, $\cW_3$, and $\cW_4$ are strictly
negative for $1\leq n\leq20$. The first coefficient is strictly
negative for every $\cW_N$, and the negativity deepens with
higher-spin defect:
$\tilde\lambda_1(\cW_N) \sim -\log N
- (1 - \gamma - \zeta'(2)/\zeta(2))$
as $N \to \infty$, with
$1 - \gamma - \zeta'(2)/\zeta(2) \approx 0.147$.

There are two distinct Li theories:
\begin{enumerate}[label=\textup{(\roman*)}]
\item
the \emph{zero-side} Li--Keiper coefficients, attached to the
Eisenstein/Rankin--Selberg interface and universal up to the
scalar zeta kernel~$\varepsilon^1_s = 4\zeta(2s)$; these do
not distinguish families;
\item
the \emph{prime-side} Li coefficients~$\tilde\lambda_n(\cA)$,
attached to the connected sewing Dirichlet
lift~$S_\cA(u)$; these are
\emph{family-sensitive} and distinguish Heisenberg, Virasoro,
and~$\cW_N$ at $n = 1$.
\end{enumerate}
The zero-side coefficients probe the pole locations
of~$\zeta(2s)/\zeta(2s{-}1)$; the prime-side coefficients
probe the Euler defect of the sewing operator.
Neither alone implies the other.
\end{remark}

\subsection{The sewing Gram matrix}%
\label{subsec:surface-moment-matrix}%
\label{subsec:sewing-gram-matrix}%
\index{Gram matrix!sewing|textbf}%
\index{moment matrix!surface|textbf}

\begin{definition}[Sewing moment pairing]%
\label{def:sewing-inner-product}%
\index{sewing inner product|textbf}%
\index{Dirichlet polynomial}%
On finite linear combinations of power functions on $\mathbb Z_{>0}$,
define
\begin{equation}\label{eq:sewing-inner-product}
\langle f, g \rangle_\cA
\;:=\;
\sum_{N=1}^\infty
a_\cA(N)\,f_N\,\overline{g_N},
\end{equation}
where $a_\cA(N)$ are the connected sewing amplitudes
(Definition~\ref{def:connected-genus1-free-energy}).
For the Heisenberg algebra, the weight is the
multiplicative divisor function:
$a_\cH(N) = \sigma_{-1}(N) = \sum_{d \mid N}d^{-1}$.
\end{definition}

\begin{definition}[Surface moment matrix]%
\label{def:surface-moment-matrix}%
\index{Hankel matrix!surface}%
For $\alpha \geq 2$, the \emph{surface moment matrix}
\begin{equation}\label{eq:surface-moment-matrix}
\mathscr{M}_\cA(\alpha)_{ij}
\;:=\; S_\cA(\alpha + i + j)
\;=\;
\langle N^{-\alpha/2-i},\, N^{-\alpha/2-j}\rangle_\cA
\end{equation}
is the Gram matrix of the power functions
$\{N^{-\alpha/2-j}\}_{j \geq 0}$
in the sewing moment pairing~\eqref{eq:sewing-inner-product}.
Write $\mu_k(\cA;\alpha) :=
\det\bigl(\mathscr{M}_\cA(\alpha)|_{[0,k)^2}\bigr)$.
\end{definition}

\begin{theorem}[Gram positivity; \ClaimStatusProvedHere]%
\label{thm:surface-moment-positivity}%
\label{thm:gram-positivity}%
\index{moment matrix!positivity}%
\index{Gram matrix!positivity}%
$\mathscr{M}_\cA(\alpha) \succ 0$
for every $\alpha \geq 2$ and every chiral algebra~$\cA$
whose sewing coefficients satisfy $a_\cA(N) \geq 0$ and have infinite
support.
In particular, $\mu_k(\cA;\alpha) > 0$ for all $k \geq 1$.
\end{theorem}

\begin{proof}
$\mathscr{M}_{ij} = \sum_N a_\cA(N)\,(N^{-\alpha/2-i})(N^{-\alpha/2-j})
= V^T\!DV$,
where $D = \operatorname{diag}(a_\cA(N))_{N \geq 1} \succeq 0$
and $V_{Nj} = N^{-\alpha/2-j}$.
Every finite principal minor is a weighted Vandermonde Gram matrix on
the distinct bases $N^{-1}$. Since the support of $a_\cA$ is infinite,
one may choose at least $k$ positive-weight bases for the $k$-th minor,
and the corresponding Vandermonde determinant is nonzero.
\end{proof}

\begin{corollary}[Virasoro sewing deficit; \ClaimStatusProvedHere]%
\label{cor:virasoro-gram-ratio}%
\index{Gram matrix!Virasoro ratio}%
\begin{equation}\label{eq:surface-minor-ratio}
\frac{\mathscr M_{\mathrm{Vir}}(\alpha)_{ij}}
{\mathscr M_{\cH}(\alpha)_{ij}}
=
1-\zeta(\alpha+i+j)^{-1}
\qquad(\alpha\geq2,\ i,j\geq0).
\end{equation}
Equivalently,
$a_{\mathrm{Vir}}(N)=\sigma_{-1}(N)-N^{-1}$: the Virasoro sewing
measure is obtained from the Heisenberg divisor measure by removing
the missing weight-$1$ oscillator mode. No determinant ratio for
$k\geq2$ follows from the entrywise identity alone.
\end{corollary}

\begin{proof}
$S_{\mathrm{Vir}}(u) = (1 - \zeta(u)^{-1})\cdot S_\cH(u)$
gives the entrywise identity after substituting
$u=\alpha+i+j$. Since
$S_\cH(u)=\sum_N\sigma_{-1}(N)N^{-u}$ and
$S_{\mathrm{Vir}}(u)=\zeta(u{+}1)(\zeta(u)-1)$, the coefficient
identity is
$a_{\mathrm{Vir}}(N)=\sigma_{-1}(N)-N^{-1}$.
\end{proof}

\begin{definition}[Sewing reproducing kernel]%
\label{def:sewing-kernel}%
\index{reproducing kernel!sewing|textbf}%
The \emph{sewing reproducing kernel} is the bilinear Dirichlet kernel
\begin{equation}\label{eq:sewing-kernel}
K_\cA(s,t) \;:=\; S_\cA(s+t)
\;=\; \langle N^{-s},\, N^{-t}\rangle_\cA
\end{equation}
on $\{(s,t) : \operatorname{Re}(s+t) > 1\}$.
The associated Hermitian RKHS kernel is
$K_\cA^{\mathrm{Hilb}}(s,t):=K_\cA(s,\bar t)=S_\cA(s+\bar t)$.
For the Heisenberg algebra:
$K_\cH(s,t) = \zeta(s{+}t)\,\zeta(s{+}t{+}1)$.
\end{definition}

\begin{theorem}[Sewing RKHS; \ClaimStatusProvedHere]%
\label{thm:sewing-rkhs}%
\index{reproducing kernel Hilbert space!sewing|textbf}%
Assume $S_\cA(u)=\sum_N a_\cA(N)N^{-u}$ converges absolutely for
$\operatorname{Re}(u)>1$ and that $a_\cA(N)>0$ on its support.
The Hermitian sewing kernel $K_\cA^{\mathrm{Hilb}}$ generates a
reproducing kernel Hilbert space $\cH_\cA$ of Dirichlet
series $f(s) = \sum_N c_N\,N^{-s}$ with norm
\begin{equation}\label{eq:sewing-rkhs-norm}
\|f\|_\cA^2
\;=\;
\sum_{\substack{N\geq 1\\a_\cA(N)>0}}
\frac{|c_N|^2}{a_\cA(N)}\,,
\end{equation}
convergent on $\operatorname{Re}(s) > 1$.
The reproducing property is
$f(s) = \langle f,\,
K_\cA^{\mathrm{Hilb}}(\cdot, s)\rangle_{\cH_\cA}$.
For Heisenberg: $\|f\|_\cH^2 = \sum_N |c_N|^2/\sigma_{-1}(N)$,
a weighted Hardy space with multiplicative divisor weight.
\end{theorem}

\begin{proof}
$K_\cA^{\mathrm{Hilb}}(s,t)
= \sum_N a_\cA(N)\,N^{-s}\,N^{-\bar{t}}$
is a positive-definite kernel by
Theorem~\ref{thm:gram-positivity}.
By the Moore--Aronszajn theorem, it generates a unique
RKHS. The norm~\eqref{eq:sewing-rkhs-norm} follows from
the kernel expansion: the functions
$e_N(s) := \sqrt{a_\cA(N)}\,N^{-s}$ form an orthonormal
basis, and
$K_\cA^{\mathrm{Hilb}}(s,t)
= \sum_N e_N(s)\,\overline{e_N(t)}$.
A general element $f = \sum c_N N^{-s}$ has
$\|f\|^2 = \sum |c_N|^2/a_\cA(N)$.
\end{proof}

\begin{construction}[Interacting sewing kernel]%
\label{constr:interacting-kernel}%
\index{reproducing kernel!interacting|textbf}%
\index{shadow obstruction tower!kernel corrections}%
The shadow obstruction tower
(Definition~\ref{def:shadow-postnikov-tower})
provides corrections to the free-field kernel.
Define the \emph{interacting sewing kernel} at degree
cutoff~$R$:
\begin{equation}\label{eq:interacting-kernel}
K_\cA^{\leq R}(s,t)
\;:=\;
K_\cA(s,t)
+ \sum_{r=3}^{R}
\operatorname{Sh}_r\bigl(\Theta_\cA^{\leq r};\,s,t\bigr),
\end{equation}
where the degree-$r$ shadow kernel is
\begin{equation}\label{eq:shadow-kernel-mellin}
\operatorname{Sh}_r(\Theta_\cA^{\leq r};\,s,t)
\;:=\;
\sum_{N=1}^\infty b_r(N)\,N^{-s-t},
\quad
b_r(N) := \sum_{\substack{\Gamma:\,\text{genus-1}\\
|V(\Gamma)|=r}}
\frac{\ell_\Gamma(N)}{|\operatorname{Aut}(\Gamma)|},
\end{equation}
with $\ell_\Gamma(N)$ the $N$-th Fourier coefficient of
the graph amplitude on the elliptic curve
(Theorem~\ref{thm:sewing-shadow-intertwining}).
At $R = 2$: $b_2(N) = 0$ for all $N$, so
$K_\cA^{\leq 2} = K_\cA$.
At $R = 4$: $b_4$ carries the quartic resonance class
$\mathfrak{Q}(\cA)$.
\end{construction}

\begin{proposition}[Collapse permanence; \ClaimStatusConditional]%
\label{prop:collapse-permanence}%
\index{sewing--Hecke reciprocity!collapse permanence}%
\index{interacting sewing kernel!factorisation}%
For every degree cutoff~$R$, under the Mellin scalar projection used
in Construction~\ref{constr:interacting-kernel}, the interacting
sewing kernel
$K_\cA^{\leq R}(s,t)$ factors through $s{+}t$:
\begin{equation}\label{eq:collapse-permanence}
K_\cA^{\leq R}(s,t)
\;=\;
\sum_{N=1}^\infty
\Bigl(a_\cA(N) + \sum_{r=3}^{R}b_r(N)\Bigr)\,N^{-s-t}.
\end{equation}
Consequently, the Sewing--Hecke collapse
\textup{(}Theorem~\textup{\ref{thm:sewing-hecke-reciprocity}}\textup{)}
is permanent: shadow-tower interactions modify the Dirichlet
coefficients of~$S_\cA$ but do not break the collapse
to a single variable $v = s{+}u{-}1$.
\end{proposition}

\begin{proof}
The graph-amplitude Fourier coefficients
$b_r(N) = \sum_\Gamma \ell_\Gamma(N)/|\operatorname{Aut}(\Gamma)|$
are scalars depending only on~$N$, not on $(s,t)$.
The Dirichlet series~\eqref{eq:shadow-kernel-mellin}
is $\sum_N b_r(N)\,N^{-s-t}$, which factors through
$s{+}t$. Summing over~$r$ preserves the factorisation.
The content is the choice of scalar Mellin projection: the graph data
enter as Fourier coefficients $b_r(N)$. Non-scalar shadow tensors may
retain more variables before this projection.
The Rankin--Selberg unfolding argument of
Theorem~\ref{thm:sewing-hecke-reciprocity} applies
term-by-term, giving the interacting two-variable
$L$-object $L_\cA^{\leq R}(s,u) =
\Gamma(v)\,(2\pi)^{-v}\,S_\cA^{\leq R}(v)$
with $S_\cA^{\leq R}(v) := \sum_N
(a_\cA(N) + \sum_{r=3}^R b_r(N))\,N^{-v}$.
\end{proof}

\begin{proposition}[\ClaimStatusProvedHere]%
\label{prop:benjamin-chang-bridge}%
\index{Benjamin--Chang}%
\index{constrained Epstein zeta}%
Let $\varepsilon^c_s = \sum_{\Delta \in S}(2\Delta)^{-s}$
denote the constrained Epstein zeta of
Benjamin--Chang~\cite{BenjaminChang22}.
At the self-dual radius $R = 1$:
\begin{equation}\label{eq:bc-sewing-bridge}
\varepsilon^1_s(R{=}1)
= 4\,\zeta(2s)
= \frac{4}{\zeta(2s{+}1)}\cdot K_\cH(s,s).
\end{equation}
The diagonal of the sewing kernel $K_\cH(s,s) = S_\cH(2s)$
recovers the primary-counting Dirichlet series up to the
factor $4/\zeta(2s{+}1)$, which is the plethystic ratio
\textup{(}connected vs.\ disconnected\textup{)}.
\end{proposition}

\begin{proof}
$\varepsilon^1_s(R{=}1) = 4\zeta(2s)$
(self-dual scalar primaries $\Delta = k^2/2$,
multiplicity~$4$).
\end{proof}

\begin{theorem}[\ClaimStatusConditional]%
\label{thm:shadow-euler-independence}%
\index{shadow depth!vs Euler--Koszul}%
\index{Euler--Koszul!vs shadow depth}%
Shadow depth $\kappa_d$
\textup{(}Definition~\textup{\ref{def:shadow-depth-classification}}\textup{)}
and Euler--Koszul class $\operatorname{ek}$
\textup{(}Definition~\textup{\ref{def:euler-koszul-class}}\textup{)}
are independent:
\begin{center}
\small
\renewcommand{\arraystretch}{1.2}
\begin{tabular}{lcccc}
\toprule
\emph{Family} & $W(\cA)$ & $\kappa_d$ & $\operatorname{ek}$
 & $K_\cA(s,t)$ \\
\midrule
$\cH_c$ & $\{1\}$ & $2$ & $0$
 & $\zeta(s{+}t)\,\zeta(s{+}t{+}1)$ \\
$V_k(\fg)$ & $\{1^{\dim\fg}\}$ & $3$ & $0$
 & $\dim\fg\cdot\zeta\,\zeta$ \\
$\beta\gamma$ & $\{1,1\}$ & $4$ & $0$
 & $2\,\zeta\,\zeta$ \\
$\mathrm{Vir}_c$ & $\{2\}$ & $\infty$ & $1$
 & $\zeta(s{+}t{+}1)(\zeta(s{+}t){-}1)$ \\
$\cW_N$ & $\{2,\ldots,N\}$ & $\infty$ & $N{-}1$
 & \emph{harmonic defect} \\
\bottomrule
\end{tabular}
\end{center}
The Gram matrix $\mathscr{M}_\cA(\alpha) \succ 0$ for all
five families \textup{(}Theorem~\textup{\ref{thm:gram-positivity}}\textup{)}.
The bare kernel $K_\cA$ detects the Euler--Koszul defect. It is not a
shadow-depth detector by itself: exact Euler--Koszul families already
split into $\cH$, $V_k(\fg)$, and $\beta\gamma$ with
$\kappa_d=2,3,4$. Interacting corrections vanish precisely on the
Gaussian depth-$2$ lane; for $\beta\gamma$ the quartic shadow gives a
nonzero $b_4$ even though $\operatorname{ek}=0$.
\end{theorem}

\begin{proof}
Three witnesses:
$\cH$, $V_k(\fg)$, $\beta\gamma$ share $\operatorname{ek} = 0$
with $\kappa_d \in \{2,3,4\}$;
$\mathrm{Vir}$ and $\cW_3$ share $\kappa_d = \infty$ with
$\operatorname{ek} \in \{1, 2\}$.
Shadow depths: Theorem~\ref{thm:shadow-archetype-classification}.
Euler--Koszul classes: Theorem~\ref{thm:euler-koszul-standard}.
For exact Euler--Koszul families, $D_\cA\equiv1$, so
$K_\cA=\operatorname{rk}(\cA)\cdot\zeta\,\zeta$ has no arithmetic
defect. This does not force $K_\cA^{\leq R}=K_\cA$:
$\beta\gamma$ has $\operatorname{ek}=0$ but class~C shadow depth~$4$,
so its quartic shadow contributes at $R=4$. The interaction-free
condition is $\kappa_d=2$, not $\operatorname{ek}=0$.
\end{proof}

\begin{corollary}[\ClaimStatusConditional]%
\label{cor:two-faces-theta}%
\index{Theta@$\Theta_\cA$!two faces}%
$\operatorname{ek}(\cA)$ depends only on the graded character, while
$\kappa_d(\cA)$ depends on the OPE. The bare sewing reproducing kernel
$K_\cA(s,t)=S_\cA(s+t)$ sees the Euler--Koszul defect but not shadow
depth. The pair
\[
\bigl(K_\cA,\{b_r\}_{r\geq3}\bigr),
\]
or equivalently the finite-degree projections
$\Theta_\cA^{\leq r}$, records both faces.
\end{corollary}

\begin{proof}
$\operatorname{ek} = \max_i(w_i - 1)$: graded character.
$\kappa_d$: $A_\infty$-brackets
(Theorem~\ref{thm:shadow-archetype-classification}).
Witness: $\beta\gamma$ and $\cH \oplus \cH$ share
$\operatorname{ek} = 0$ and character $\eta(q)^{-2}$
but differ in $\kappa_d$ ($4$ vs.\ $2$). Since both have the same
bare Euler product, $K_\cA$ cannot distinguish them; the quartic
coefficient $b_4$ does.
\end{proof}

\begin{definition}[Genus-$1$ modular datum]%
\label{def:platonic-arithmetic-package}%
\label{def:genus1-modular-datum}%
\index{genus-1 modular datum|textbf}%
\textbf{Type signature.} (Open quadrant; level~$5$, modular shadow on
the genus-$1$ open factorization category; hypothesis package: $\cA$
modular Koszul, finite-type or weight-completed bar coalgebra,
HS-sewing where partition-function statements are made.)
The \emph{genus-$1$ modular datum} of~$\cA$ is the sextuple
\begin{equation}\label{eq:full-genus1-package}
\mathscr{D}_1(\cA)
\;:=\;
\bigl(
\kappa(\cA),\;
\Delta_\cA,\;
K_\cA,\;
\mathfrak{Q}(\cA),\;
\kappa_d(\cA),\;
\operatorname{ek}(\cA)
\bigr),
\end{equation}
encoding: the scalar genus-$1$ curvature class
\textup{(}Theorem~D\textup{)}, the spectral discriminant
\textup{(}Definition~\textup{\ref{def:shadow-postnikov-tower}}\textup{)},
the sewing reproducing kernel
\textup{(}Definition~\textup{\ref{def:sewing-kernel}}\textup{)},
the quartic resonance class
\textup{(}Appendix~\textup{\ref{app:nonlinear-modular-shadows}}\textup{)},
the shadow depth
\textup{(}Definition~\textup{\ref{def:shadow-depth-classification}}\textup{)},
and the Euler--Koszul class
\textup{(}Definition~\textup{\ref{def:euler-koszul-class}}\textup{)}.
The Gram matrix $\mathscr{M}_\cA(\alpha) \succ 0$
\textup{(}Theorem~\textup{\ref{thm:gram-positivity}}\textup{)}
is the positive-definiteness witness for $K_\cA$.
The label \emph{modular} refers to the level-$5$ trace + clutching
shadow on $\overline{\cM}_{1,n}$; modular invariance of any
partition-function specialization is the modular reconstruction of
the open category, not an intrinsic property of $\cA$ as a closed
algebra.
\end{definition}

\subsection{The Euler--Koszul classification}%
\label{subsec:euler-koszul-classification}%
\index{Euler--Koszul!classification|textbf}

\begin{definition}[Euler--Koszul tier]%
\label{def:euler-koszul-tier}%
\index{Euler--Koszul!tier|textbf}%
A modular Koszul algebra~$\cA$ is:
\begin{enumerate}[label=\textup{(\roman*)}]
\item
\emph{exact Euler--Koszul} if $D_\cA(u) \equiv 1$, equivalently
$\operatorname{ek}(\cA) = 0$;
\item
\emph{finitely defective Euler--Koszul} if
$D_\cA(u) \not\equiv 1$ but $1 - D_\cA(u)$
is a ratio of finite harmonic sums $H_n(u)$ and~$\zeta(u)$:
the defect is a \emph{finite} subtraction from the
pure Euler product;
\item
\emph{non-Euler--Koszul} if $S_\cA$ decomposes as a linear
combination of Hecke $L$-functions attached to distinct
Gr\"o\ss encharacters, so that no single Euler product
suffices.
\end{enumerate}
\end{definition}

\begin{theorem}[Euler--Koszul tier placement for standard examples;
\ClaimStatusConditional]%
\label{thm:euler-koszul-tier-classification}%
\index{Euler--Koszul!tier classification}%
\begin{enumerate}[label=\textup{(\roman*)}]
\item
Heisenberg, $V_k(\fg)$, and $\beta\gamma$ are exact
Euler--Koszul: $S_\cA = |W|\cdot\zeta(u)\,\zeta(u{+}1)$.
\item
$\mathrm{Vir}_c$ and principal $\cW_N$ are finitely defective
Euler--Koszul: the defect
$1 - D_\cA(u)$ is a finite harmonic subtraction.
\item
Positive binary quadratic Epstein zeta functions of discriminant
$D$ with class number $h(D) \geq 2$ are non-Euler--Koszul: their
genus decomposes into linear combinations of inequivalent Hecke
$L$-functions.
\end{enumerate}
\end{theorem}

\begin{proof}
(i)~Theorem~\ref{thm:euler-koszul-standard}(i).
(ii)~For $\cW_N$: $D_{\cW_N}(u) = 1 -
\sum_{j=1}^{N-1}H_j(u)/((N{-}1)\zeta(u))$, and
each $H_j$ is a finite sum.
(iii)~For a positive binary quadratic form with $h(D) \geq 2$, the
Epstein zeta function
$\varepsilon^c_s$
decomposes into a linear combination of Hecke
$L$-functions attached to distinct Gr\"o\ss encharacters;
no single Euler product suffices
(Davenport--Heilbronn).
\end{proof}

\begin{remark}[Drinfeld--Sokolov reduction as arithmetic defect]%
\label{rem:ds-arithmetic-defect}%
\index{Drinfeld--Sokolov!arithmetic defect}%
Principal DS reduction $V_k(\fg) \leadsto \cW_k(\fg)$
sends exact Euler--Koszul to finitely defective
Euler--Koszul.
Arithmetically, this is visible as:
\[
\text{pure Heisenberg Euler product}
\;\longmapsto\;
\text{Euler product} - \text{finite harmonic defect}.
\]
The defect $P_{\cW_N}(u) =
\sum_{j=1}^{N-1}H_j(u)/(N{-}1)$ encodes the higher-spin
subtraction: each generator of weight $s \geq 2$ removes
the modes $m = 1,\ldots,s{-}1$ from the Euler product,
contributing $H_{s-1}(u)/\zeta(u)$ to the defect.
\end{remark}

\subsection{Sewing--Hecke reciprocity}%
\label{subsec:two-variable-L}%
\label{subsec:sewing-hecke-reciprocity}%
\index{sewing--Hecke reciprocity|textbf}%
\index{L-object@$L$-object!two-variable|textbf}

\begin{definition}[Two-variable $L$-object]%
\label{def:two-variable-L}%
\index{L-object@$L$-object!definition|textbf}%
Let $\cA$ be a modular Koszul algebra,
$E^*(\tau,s) := \pi^{-s}\Gamma(s)\sum_{(m,n)\neq(0,0)}
y^s/|m\tau{+}n|^{2s}$
the completed Eisenstein series, and
$F_\cA^{\mathrm{conn}}(\tau)$
the connected genus-$1$ free energy
(Definition~\ref{def:connected-genus1-free-energy}).
The \emph{two-variable $L$-object} is
\begin{equation}\label{eq:two-variable-L}
L_\cA(s,u)
\;:=\;
\int_{\cF}^{\mathrm{reg}}
F_\cA^{\mathrm{conn}}(\tau)\,
E^*(\tau,s)\,
y^{u-2}\,dx\,dy,
\end{equation}
where $\cF = \mathrm{SL}(2,\bZ)\backslash\cH$ and the
regularization is Zagier truncation
$\cF^T := \{\tau \in \cF : \mathrm{Im}(\tau) \leq T\}$,
$T \to \infty$.
\end{definition}

\begin{theorem}[Diagonal Sewing--Hecke reciprocity]%
\label{thm:sewing-hecke-reciprocity}%
\ClaimStatusConditional%
\index{sewing--Hecke reciprocity!theorem|textbf}%
Assume the connected genus-$1$ sewing free energy is the diagonal
weight-multiset series of
Definition~\ref{def:weight-multiset}.  Then the two-variable
$L$-object collapses to a one-variable object:
\begin{equation}\label{eq:sewing-hecke}
L_\cA(s,u)
\;=\;
\frac{\Gamma(v)}{(2\pi)^v}\,S_\cA(v),
\qquad
v := s + u - 1,
\end{equation}
where $S_\cA$ is the diagonal Dirichlet--sewing lift
\textup{(}Theorem~\textup{\ref{thm:dirichlet-weight-formula}}\textup{)}.
The $(s,u)$-dependence factors through~$v$ exactly.
In particular\textup:
\begin{enumerate}[label=\textup{(\roman*)}]
\item
\textup{(}Sewing--Selberg recovery.\textup{)}
$L_\cA(s,0) = \Gamma(s{-}1)\,(2\pi)^{-(s-1)}\,S_\cA(s{-}1)$.
For Heisenberg, this recovers the regularized
sewing--Selberg formula
\textup{(}Theorem~\textup{\ref{thm:sewing-selberg-formula}}\textup{)}.
\item
\textup{(}Dirichlet recovery.\textup{)}
$L_\cA(1,u) = \Gamma(u)\,(2\pi)^{-u}\,S_\cA(u)$,
recovering the Mellin representation of the
Dirichlet--sewing lift.
\item
\textup{(}Zeta-zero locus.\textup{)}
The factor $\zeta(v{+}1)$ in $S_\cA(v)$ has \emph{zeros}
at $v = \rho - 1$ for each nontrivial zero~$\rho$
of~$\zeta$, tracing the line $s + u = \rho$
in the $(s,u)$-plane. These are zeros of~$L_\cA$,
not poles; the scattering poles at $s = \rho/2$ are
holomorphic points of the Rankin--Selberg integral
\textup{(}Proposition~\textup{\ref{prop:scattering-residue}}\textup{)}.
\item
\textup{(}Heisenberg.\textup{)}
$L_\cH(s,u) = \Gamma(v)\,(2\pi)^{-v}\,
\zeta(v)\,\zeta(v{+}1)$.
\item
\textup{(}Virasoro.\textup{)}
$L_{\mathrm{Vir}}(s,u) = \Gamma(v)\,(2\pi)^{-v}\,
\zeta(v{+}1)\bigl(\zeta(v) - 1\bigr)
= L_\cH(s,u)\cdot D_{\mathrm{Vir}}(v)$.
\end{enumerate}
\end{theorem}

\begin{proof}
Under the diagonal weight-multiset hypothesis, the sewing operator
$K_q$ has diagonal eigenvalues $|q|^n=e^{-2\pi ny}$ on the selected
mode basis, so $F_\cA^{\mathrm{conn}}$ depends on
$y=\mathrm{Im}(\tau)$ alone. The zeroth Fourier mode is
$\int_0^1 F_\cA^{\mathrm{conn}}(x{+}iy)\,dx
= F_\cA^{\mathrm{conn}}(e^{-2\pi y})$.
Rankin--Selberg unfolding replaces the modular integral by
\[
L_\cA(s,u)
\;=\;
\int_0^\infty
F_\cA^{\mathrm{conn}}(e^{-2\pi y})\,y^{s+u-2}\,dy.
\]
Set $v = s + u - 1$.
By the weight-multiset expansion
(Definition~\ref{def:weight-multiset}):
\begin{align*}
L_\cA(s,u)
&= \sum_{i=1}^r\,\sum_{m \ge w_i}\,\sum_{k \ge 1}
\frac{1}{k}\int_0^\infty e^{-2\pi mky}\,y^{v-1}\,dy
\\
&= \frac{\Gamma(v)}{(2\pi)^v}
\sum_{i=1}^r\,\sum_{m \ge w_i} m^{-v}\cdot\zeta(v{+}1)
\\
&= \frac{\Gamma(v)}{(2\pi)^v}\,S_\cA(v).
\qedhere
\end{align*}
\end{proof}

\begin{remark}[Interpretation of the collapse]%
\label{rem:sewing-hecke-collapse}%
\index{sewing--Hecke reciprocity!collapse}%
The collapse $L_\cA(s,u) = f(s{+}u{-}1)$ is a theorem
expressing the fact that the connected sewing free energy
is $y$-dependent alone (the sewing operator is diagonal
in the Fock basis). The zero-side variable~$s$
(Rankin--Selberg) and the prime-side variable~$u$
(Dirichlet--sewing) are \emph{locked together} by the
Mellin transform.
The regularized sewing--Selberg formula
(Theorem~\ref{thm:sewing-selberg-formula}) and the
Dirichlet--sewing lift
(Theorem~\ref{thm:dirichlet-weight-formula}) are two
\emph{faces of the same Mellin transform}, related by
the substitution $s \leftrightarrow u$ along the
constraint $v = s + u - 1$.

The zeta-zero divisors in the collapsed object are read in the
$v$-coordinate. The universal factor $\zeta(v{+}1)$ gives zeros at
$v=\rho-1$; whenever an additional $\zeta(v)$ factor survives, it gives
zeros at $v=\rho$. At the Sewing--Selberg point $u=0$ these become
lines $s=\rho$ and $s=\rho+1$, respectively. The Eisenstein scattering
poles at $s=\rho/2$ belong to the scattering coefficient, not to the
collapsed Dirichlet factor $S_\cA(v)$.

Distinct families can share the same spectral
discriminant $\Delta_\cA$ while having different sewing
lifts: $\widehat{\mathfrak{sl}}_2$, $\mathrm{Vir}_c$,
and $\beta\gamma$ all sit on $\Delta_\cA(x) = (1{-}3x)(1{+}x)$,
yet $S_{\beta\gamma} = 2\zeta(v)\zeta(v{+}1)$,
$S_{\mathrm{Vir}} = \zeta(v{+}1)(\zeta(v){-}1)$,
$S_{\widehat{\mathfrak{sl}}_2} = 3\zeta(v)\zeta(v{+}1)$.
The regularized sewing--Selberg formula and the Higgs field of
Theorem~\ref{thm:shadow-higgs-field}
are related: the sewing operator $K_q$ is the analytic avatar of the
Higgs field $\Phi_\cA$ on the shadow triple, and the
Rankin--Selberg $L$-function $L_\cA(s,u)$ encodes the spectral data
of $\theta_\cA$ along the critical lines.
\end{remark}

\section{Genus-$2$ Siegel character $(\Phi_{10}^{\mathrm{un}})^{-1}$ at $K3 \times E$}
\label{sec:genus-complete-supplement}

\begin{remark}[Genus-$2$ Siegel character $1/\Phi_{10}^{\mathrm{un}}$]
\label{rem:genus-complete-character}
\textbf{Type signature.} (Open quadrant Vol~I scalar shadow;
cross-volume bridge to the Vol~III CY-quadrant CY$_3$ stage at
$(\Sigma_{d-1}, C)=(\Sigma_2, E_\tau)$ via
$\Phi_3^{(\Sigma_2,E_\tau)}=\mathrm{Sp}^{\mathrm{ch}}_{\Sigma_2,E_\tau}\circ
\Phi_3^{\mathrm{FA}}$; level~$5$ scalar shadow on
$\overline{\cM}_2$.)

The Vol~III $K3 \times E$ Hall--Borcherds recognition target
$\mathbf H_{\Delta_5}$ has scalar denominator character
$\mathrm{Ch}_{\mathrm{sc}}(\mathbf{H}_{\Delta_5})(Z) =
1/\Phi_{10}^{\mathrm{un}}(Z)$ on the Siegel upper half-space
$\mathbb{H}_2$; at a separating-node degeneration $Z \to \mathrm{diag}(\tau, \rho)$ the
character factorises as
\[
 1/\Phi_{10}^{\mathrm{un}}(Z) \;\longrightarrow\; \frac{1}{\eta(\tau)^{24}} \cdot \frac{1}{\phi_{10,1}(\rho)},
\]
recovering the $\eta^{24}$ genus-$1$ denominator of the K3 chiral algebra times the
Jacobi form of the elliptic direction. The normalisation is fixed as follows:
$\Delta_5$ is the primitive weight-$5$ BKM denominator, the monic product is
$D_5=64^{-1}\Delta_5$, the unnormalised Igusa square is
$\Phi_{10}^{\mathrm{un}}=\Delta_5^2$, the Oberdieck--Pixton scalar is
$\Phi_{10}^{\mathrm{OP}}=D_5^2=4096^{-1}\Delta_5^2$, and the reduced
OP/DT partition scalar is
$Z_{\mathrm{OP/DT}}=-D_5^{-2}=-4096\,\Delta_5^{-2}$. The character in this
remark uses $\Phi_{10}^{\mathrm{un}}$, not the OP-normalised square. Hence the
BKM weight is
$\kappa_{\mathrm{BKM}}(\mathfrak{g}_{\Delta_5}) = c_\Lambda(0)/2 =
\mathrm{wt}(\Delta_5) = 5$, in agreement with the universal
Borcherds-weight identity
$\kappa_{\mathrm{BKM}}(\Phi_N) = c_N(0)/2$ at $N=\mathrm{wt}(\Delta_5)$.
The $\mathrm{Sp}_4(\Z)$-modularity of $\Phi_{10}^{\mathrm{un}}$ is the
level-$5$ modular reconstruction on the open factorization category at
$(\Sigma_2,E_\tau)$: the closed shadow inherits Siegel modularity, but
the modular property itself is the trace + clutching equivariance, not
an intrinsic closed-algebra invariant. Promotion from this scalar
denominator identity to a full chiral bialgebra trace requires the
Vol~III Hall--Borcherds recognition and gerbe-banded Tannakian
reconstruction hypotheses; the scalar automorphic identity itself is
the unconditional input.
\end{remark}

\begin{remark}[Bi-based factorisation and higher-genus averaging]
\label{rem:genus-complete-bi-based}
The higher-genus completion at $K3 \times E$ is bi-based: chiral base
$\mathrm{Ran}(E^{\mathrm{nod,sm}}_{24})$ on the smooth locus of the $24$-node
discriminant curve, with nearby-cycles $\mathcal{D}$-modules at each Kodaira $I_1$-node;
parameter base $\overline{\mathcal{A}_2}$ with Francis--Gaitsgory factorisation
$\infty$-category in holonomic $\mathcal{D}$-modules regular-singular on Humbert
$H_1 \cup H_4$. Averaging morphism
$\mathrm{av} = \mathrm{Torelli}_{K3} \circ \mathrm{KugaSatake} \circ j_{\mathrm{Kodaira}}
\colon \mathrm{Ran}(\mathbb{P}^1)_{M_{24}} \to \overline{\mathcal{A}_2}$. The BD chirality
axioms require smooth base, so nearby-cycles continuation at the $24$ nodes is
load-bearing.
\end{remark}

\section{Genus-$2$ Tannakian avatar}
\label{sec:genus-complete-supplement-ii}

\begin{remark}[Tannakian fibre functor at genus $2$]
\label{rem:genus-complete-tannakian}
\textbf{Type signature.} (Cross-volume Open$\leftrightarrow$CY at
level~$4$/$5$; Vol~I open factorization category at $(\Sigma_2,
E_\tau)$ versus Vol~III CY$_3$ Hall--Borcherds bulk; hypothesis
package: gerbe-banded Tannakian reconstruction on $U$, conformal OS
positivity for the scalar trace, HS-sewing on the open lane.)

The Siegel-genus-$2$ scalar character $1/\Phi_{10}^{\mathrm{un}}(Z)$ of
Remark~\ref{rem:genus-complete-character} is the
\emph{Tannakian graded-trace target} of the fibre functor
$\omega\colon \mathrm{Rep}(\mathbf H_{\Delta_5}) \to
\mathrm{sVec}_{\mathbb C}$ defined on Nakajima's Hilbert-scheme
cohomology, after gerbe banding on
$U=\overline{\mathcal A_2}\setminus(H_1\cup H_4)$. Concretely, under
that recognition hypothesis,
\[
 \mathrm{Ch}\bigl(\mathbf H_{\Delta_5}\bigr)(Z)
 \;=\;
 \mathrm{tr}_{\mathrm{s}}\!\bigl(q^{L_0}\,q'^{L_0'}\mid \omega(\mathbf H_{\Delta_5})\bigr)
 \;=\;
 1/\Phi_{10}^{\mathrm{un}}(Z),
\]
with $Z = \left(\begin{smallmatrix}\tau & z\\ z & \rho\end{smallmatrix}\right)$
and $q = e^{2\pi i\tau},\ q' = e^{2\pi i\rho},\ p = e^{2\pi i z}$. The
$Z$-dependence fibres the fibre functor over the Siegel upper
half-space $\mathbb H_2$: the Tannakian fibre at the generic point
$Z\in \mathbb H_2$ is neutral (by
Vol~III Remark~\ref{V3-rem:cydkappa-gerbe}), while at the Humbert
boundary $H_1 \cup H_4$ it is $\mu_8$-gerbe-obstructed, with
monodromy order $8$ at $H_1$ and $16$ at $H_4$.
The banding is now chain-level: the explicit \v{C}ech 2-cocycle
$F_{ij}(Z)=(\Phi_{10}^{\mathrm{un}}/\eta^{24}|_{U_i}\,/\,\Phi_{10}^{\mathrm{un}}/\eta^{24}|_{U_j})^{1/8}_{\mathrm{princ}}$
on the Igusa fundamental-domain cover of $U=\overline{\mathcal A_2}\setminus(H_1\cup H_4)$
trivialises the gerbe on $U$ (Vol~III
Theorem~\ref{V3-thm:modtr-banding-cocycle}), and the
Atkin--Lehner $w_8$-multiplier composition
$\mathrm{Mon}_{\gamma_1}=\zeta_4\cdot\zeta_2=\zeta_8$ along the
transverse loop at a smooth $H_1$-point witnesses non-extendability
through the non-split extension $1\to\mu_4\to\mu_8\to\mu_2\to 1$
of Bruinier 2002 LNM~1780 Prop.~5.1 (Vol~III
Theorem~\ref{V3-thm:qgf-humbert-non-extendability}).
The scalar diagonal part of the Feynman-transform genus tower pulls
back along $\mathrm{av}$ to the determinant line of this Tannakian
trace. At $g=2$, its first Chern class is
$\kappa_{\mathrm{BKM}}\lambda_2^{\mathrm{FP}}$ on the diagonal
scalar lane; any non-diagonal Siegel or Humbert contribution belongs
to the cross-channel term, not to the scalar trace. This is the
genus-$2$ specialization of Vol~III
Remark~\ref{V3-rem:cydkappa-fibre-functor}.
\end{remark}

\begin{remark}[Mukai-grading compatibility with modular operad]
\label{rem:genus-complete-mukai-modular}
The Mukai lattice $\Lambda_{\mathrm{Muk}}$ (rank $24$) gradings
are compatible with the modular-operad gluing on
$\overline{\mathcal{M}}_{g,n}$: at a separating-node degeneration
$Z \to \mathrm{diag}(\tau,\rho)$ the character factorises as
$1/\eta^{24}(\tau) \cdot 1/\phi_{10,1}(\rho)$, and the $24$ in
$\eta^{24}$ is the rank of $H^\star(K3,\mathbb Z)$, i.e.\ the Mukai
lattice rank. The Mukai grading direction is: grading is on the
\emph{K3 direction}, parameter variation is along the
\emph{$E$-direction}. No K\"unneth doubling occurs at
separating-node factorisation: the $\eta^{24}$ factor is one copy of
the Mukai rank, not two. The bi-based Ran-factorisation of
Remark~\ref{rem:genus-complete-bi-based} encodes this
single-directional grading: the chiral base carries the factorisation
(topological gluing of surfaces), the parameter base carries the
grading (Mukai-lattice variation); see Vol~II
Remark~\ref{V2-rem:htpo-mukai-direction} and Vol~III
Remark~\ref{V3-rem:cydkappa-mukai-grading-direction}. Primary:
Getzler--Kapranov 1998 Comp.~Math.~110 (modular operad);
Mukai 1987 Nagoya Math.~J.~81 (Mukai pairing); Gritsenko--Nikulin
1998 Amer.~J.~Math.~120 (Siegel $\Phi_{10}^{\mathrm{un}}$ factorisation at
separating node).
\end{remark}
